% ============================================================
%  第四章：效用函数（精讲版）
%  编译：LuaLaTeX（两遍）
% ============================================================
% 第 4 章 统一书稿前言包装
\newcommand{\lectureno}{第四章}
\newcommand{\lectureshort}{效用函数}
% ============================================================
%  微观经济学理论课程讲义 —— 统一课程前言（四讲共用）
%  版式参考 Nolan Miller, Notes on Microeconomic Theory (SWP/notes3 风格)
%  编译：LuaLaTeX（两遍）
%  各讲 preamble.tex 只需设定 \lectureno / \lectureshort 后 \input 本文件。
% ============================================================
\documentclass[11pt,a4paper,fontset=none,oneside,openany]{ctexbook}

% ---------- 版式 ----------
\usepackage{geometry}
\geometry{
  a4paper,
  top=2.8cm, bottom=2.8cm,
  left=2.6cm, right=2.6cm,
  headsep=0.85cm, footskip=1.25cm
}
\setlength{\headheight}{14pt}

% ---------- 字体（全部矢量嵌入，保证 PDF 高清晰度）----------
% 注意：本机 macOS 已移除 Kaiti SC，CJK 斜体统一回退到 Songti SC；
%        Songti SC / Heiti SC 经 OSFONTDIR 指向 Supplemental 后由 luaotfload 解析。
\usepackage{fontspec}
\setmainfont{TeX Gyre Termes}[Ligatures=TeX]
\setsansfont{TeX Gyre Heros}[Ligatures=TeX,Scale=MatchLowercase]
\setmonofont{TeX Gyre Cursor}[Scale=MatchLowercase]
\setCJKmainfont{Songti SC}[BoldFont={Songti SC Bold},ItalicFont={Songti SC}]
\setCJKsansfont{Heiti SC}[BoldFont={Heiti SC}]

% ---------- 数学 ----------
\usepackage{amsmath}
\usepackage{mathtools}
\usepackage{amsthm}
\usepackage{unicode-math}
\setmathfont{TeX Gyre Termes Math}
\allowdisplaybreaks[2]
\setlength{\abovedisplayskip}{9pt plus 3pt minus 3pt}
\setlength{\belowdisplayskip}{9pt plus 3pt minus 3pt}

% ---------- 定理类环境（书稿格式：按章编号，如“定理 1.1”）----------
\theoremstyle{plain}
\newtheorem{theorem}{定理}[chapter]
\newtheorem{proposition}[theorem]{命题}
\newtheorem{lemma}[theorem]{引理}
\newtheorem{corollary}[theorem]{推论}
\theoremstyle{definition}
\newtheorem{definition}[theorem]{定义}
\newtheorem{condition}[theorem]{条件}
\newtheorem{example}[theorem]{例}
\theoremstyle{remark}
\newtheorem{remark}[theorem]{注}
% 公式、图、表按章编号（书稿格式：式 (1.1)、图 1.1、表 1.1）
\numberwithin{equation}{chapter}
\numberwithin{figure}{chapter}
\numberwithin{table}{chapter}

% 定理类环境标题后的标点：amsthm 默认加英文 "."，此处统一改为中文句号。
% 原理：amsthm 的 \@thm 每次开环境都执行 \thm@headpunct{.}，而样式宏在其之后展开，
% 故把 \thm@headpunct{。} 追加到三个样式宏末尾即可覆盖（用 \AtBeginDocument 以免被 ctex 覆盖）。
\makeatletter
\AtBeginDocument{%
  \let\WBth@plain\th@plain
  \let\WBth@definition\th@definition
  \let\WBth@remark\th@remark
  \def\th@plain{\WBth@plain\thm@headpunct{。}}%
  \def\th@definition{\WBth@definition\thm@headpunct{。}}%
  \def\th@remark{\WBth@remark\thm@headpunct{。}}%
}
\makeatother

% 无编号小标题（仿 notes3 的 “A Note on ...”）
\newcommand{\notehead}[1]{\par\medskip\noindent{\bfseries #1}\par\smallskip}

% ---------- 表格 ----------
\usepackage{booktabs}
\usepackage{array}
\renewcommand{\arraystretch}{1.25}
\setlength{\tabcolsep}{8pt}

% ---------- 图形 ----------
\usepackage{graphicx}
\usepackage{tikz}
\usepackage{pgfplots}
\pgfplotsset{compat=1.18}
\usetikzlibrary{arrows.meta,positioning,shapes.geometric,calc,patterns,decorations.pathreplacing,decorations.markings}
\usepackage{float}
\usepackage{caption}
\captionsetup{font=small,labelfont=bf,labelsep=quad,skip=4pt}

% 统一图形配色（黑白模式：黑/灰梯度，供所有 TikZ 图复用）
% —— 全局黑白化：原彩色调色板统一改为黑/灰，保证打印与投影均可读 ——
\definecolor{AxisBlue}{RGB}{0,0,0}        % 轴线 / 主曲线：黑
\definecolor{SoftBlue}{RGB}{130,130,130}  % 节点填充：中灰
\definecolor{LightRed}{RGB}{224,224,224}  % 浅填充：浅灰
\definecolor{DeepBlue}{RGB}{0,0,0}        % 描边：黑
\definecolor{MidBlue}{RGB}{70,70,70}      % 次曲线 / 网格：深灰
\tikzset{
  >=Stealth,
  every picture/.append style={font=\small},
  axisline/.style={thick,color=black,-Stealth},
  axislabel/.style={color=black},
  fillsoft/.style={fill=LightRed,fill opacity=0.6,draw=black,thick},
  nodeblue/.style={circle,fill=black!85,draw=black,thick,inner sep=2pt},
  nodewhite/.style={circle,fill=white,draw=black,thick,inner sep=2pt},
  arrblue/.style={->, thick, color=black},
  boxsoft/.style={draw=black, rounded corners=2pt, align=center, inner sep=5pt, fill=black!6},
  % 多曲线 / 多集合区分：以线型 + 标记代替颜色
  curveA/.style={thick,black,solid},
  curveB/.style={thick,black,dashed},
  curveC/.style={thick,black,dotted},
  markA/.style={mark=*,mark size=2.2pt,color=black},
  markB/.style={mark=square*,mark size=2.0pt,color=black},
  markC/.style={mark=triangle*,mark size=2.4pt,color=black},
  % 图内区域 / 点标签（白底避免与图形重叠）
  figlabel/.style={font=\scriptsize,fill=white,inner sep=1pt,outer sep=1pt,text=black},
}
% 图形说明（黑白模式下对图形区域 / 曲线 / 关键点的解释，置于图形下方）
\newcommand{\figexplan}[1]{%
  \par\smallskip\noindent{\small\itshape 图形说明：#1}\par\smallskip}

% ---------- 页眉页脚（书稿格式：左=书名，右=章次与章名）----------
\usepackage{fancyhdr}
% 章标题：居中、大号（书稿章首样式）
\CTEXsetup[format={\centering\LARGE\bfseries}]{chapter}
\CTEXsetup[beforeskip=16pt,afterskip=14pt]{chapter}
\CTEXsetup[format={\large\bfseries}]{section}
\CTEXsetup[format={\normalsize\bfseries}]{subsection}
\pagestyle{fancy}
\fancyhf{}
\providecommand{\booktitle}{微观经济学理论：可行集、偏好与效用}
\providecommand{\lectureno}{第一章}
\providecommand{\lectureshort}{课程讲义}
\fancyhead[L]{\small\itshape \booktitle}
\fancyhead[R]{\small \lectureno\quad \lectureshort}
\fancyfoot[C]{\small \thepage}
\renewcommand{\headrulewidth}{0.4pt}
\fancypagestyle{plain}{\fancyhf{}\fancyfoot[C]{\small \thepage}\renewcommand{\headrulewidth}{0pt}}

% ---------- 正文细节 ----------
\usepackage{setspace}
\setstretch{1.12}
\usepackage{microtype}
\tolerance=900
\emergencystretch=2em
\raggedbottom

% ---------- 超链接与 PDF 书签 ----------
\usepackage{hyperref}
\hypersetup{
  colorlinks=true,
  linkcolor=black,
  urlcolor=black,
  citecolor=black,
  pdfencoding=auto,
  bookmarks=true,
  pdftitle={\booktitle\quad \lectureno：\lectureshort},
  pdfsubject={微观经济学理论书稿}
}
% 章节标题中常见数学符号的安全书签映射（避免 "Improper alphabetic constant"）
\pdfstringdefDisableCommands{%
  \def\leftrightarrow{<->}%
  \def\star{*}%
  \def\S{S}%
  \def\succeq{\text{≽}}%
  \def\succ{\text{≻}}%
  \def\sim{\text{∼}}%
}

% ---------- 常用记号 ----------
\newcommand{\R}{\mathbb{R}}
\newcommand{\N}{\mathbb{N}}
\newcommand{\Z}{\mathbb{Z}}
\newcommand{\Q}{\mathbb{Q}}
\newcommand{\succeqq}{\succeq}
\newcommand{\argmax}{\operatorname*{arg\,max}}
\newcommand{\argmin}{\operatorname*{arg\,min}}
\newcommand{\st}{\text{s.t.}\ }

\setcounter{chapter}{3}


\begin{document}

\chapter{效用函数：函数形式大全、顶刊前沿、选择方法与政策应用}

\noindent\textbf{本章任务。}第三章回答了“偏好何时能被效用函数表示”；
本章回答下一个实务问题：\textbf{做研究时，面对具体问题，该写哪一个效用函数？}
分四个层次展开：
\begin{enumerate}
  \item \textbf{函数形式库}：静态（含需求系统）、跨期、风险与模糊、行为与自控、
        社会偏好与社会福利、离散选择六大类共 40 余种效用函数，
        逐一给出性质、推导、图形、适用情形（§4.1--§4.6）；
  \item \textbf{顶刊前沿}：2021--2026 年 AER/JPE 等顶刊如何使用效用函数、
        新在哪里（公式 + 假设 + 识别策略拆解，§4.7）；
  \item \textbf{选择方法论}：研究中挑选效用函数的十步流程与决策树（§4.8）；
  \item \textbf{政策应用}：以中国延迟退休与生育政策为例，
        完整走一遍“选函数—推导—计算—敏感性”的全流程（§4.9）。
\end{enumerate}

\noindent 所有需求函数、弹性、比较静态均给出\textbf{逐步推导}；
§4.6 的全部数值算例可以手工复算验证。

\tableofcontents
\clearpage

% ============================================================
\section{静态确定性效用函数：形式库与推导}
% ============================================================

设消费束 $x=(x_1,x_2)\in\R^2_+$（推广到 $n$ 商品处另注），价格 $p=(p_1,p_2)$，收入 $m$。
每种函数按统一格式处理：\textbf{形式 → 关键性质（含推导）→ 图形 → 适用情形}。

\subsection{总览：一条“替代弹性”谱系}

多数静态函数可以排在替代弹性 $\sigma$ 的一条谱系上：

\begin{center}
\begin{tikzpicture}[scale=1.0]
  \draw[axisline] (0,0) -- (12.8,0) node[below right, axislabel] {替代弹性 $\sigma$};
  \foreach \x/\lab in {1/{0}, 5.4/{1}, 11.8/{$\infty$}}{
    \draw[AxisBlue] (\x,0.1) -- (\x,-0.1) node[below] {\lab};
  }
  \node[nodeblue, label={above=4pt, align=center:{完全互补\\ （固定比例）}}] at (1,0) {};
  \node[nodeblue, label={above=4pt, align=center:{CES\\ $0<\sigma<1$}}] at (3.2,0) {};
  \node[nodeblue, label={above=4pt, align=center:{乘积型\\ （对数可加）}}] at (5.4,0) {};
  \node[nodeblue, label={above=4pt, align=center:{CES\\ $\sigma>1$}}] at (8.4,0) {};
  \node[nodeblue, label={above=4pt, align=center:{完全替代\\ （线性）}}] at (11.8,0) {};
  \node[align=center] at (6.4,-1.3) {拟线性、基本需求型（平移乘积型）不在此轴上：它们变动的是\textbf{收入效应}而非替代弹性};
\end{tikzpicture}
\end{center}
\figexplan{横轴为替代弹性 $\sigma$。自左至右：$\sigma=0$ 为完全互补（固定比例），$0<\sigma<1$ 为互补性较强的 CES，$\sigma=1$ 为乘积型（对数可加），$\sigma>1$ 为替代性较强的 CES，$\sigma=\infty$ 为完全替代（线性）。图下方注明：拟线性与基本需求型（平移乘积型）不在这条轴上——它们变动的是收入效应而非替代弹性。}

\subsection{完全替代（线性效用）}

\notehead{形式}
$u(x_1,x_2)=ax_1+bx_2$（$a,b>0$）。

\notehead{性质与推导}
无差异曲线是直线 $x_2=u/b-(a/b)x_1$，边际替代率恒为
\[
\mathrm{MRS}=\frac{\partial u/\partial x_1}{\partial u/\partial x_2}=\frac{a}{b}.
\]
替代弹性 $\sigma=\infty$：价格比稍有偏离，需求就完全倒向一种商品。
效用最大化 $\max ax_1+bx_2\ \st p_1x_1+p_2x_2=m$：
比较“每单位收入的边际效用” $a/p_1$ 与 $b/p_2$，
\[
x_1^*=
\begin{cases}
m/p_1, & a/p_1>b/p_2,\\[2pt]
[0,\,m/p_1]\ \text{任意}, & a/p_1=b/p_2,\\[2pt]
0, & a/p_1<b/p_2,
\end{cases}
\qquad
v(p,m)=m\cdot\max\Bigl\{\frac{a}{p_1},\frac{b}{p_2}\Bigr\}.
\]
支出函数 $e(p,u)=u\cdot\min\{p_1/a,\ p_2/b\}$。

\begin{center}
\begin{tikzpicture}[scale=0.9]
  \draw[axisline] (0,0) -- (5.4,0) node[below left, axislabel] {$x_1$};
  \draw[axisline] (0,0) -- (0,4.4) node[below left, axislabel] {$x_2$};
  % 无差异曲线：平行直线族，斜率 -a/b = -0.9（比预算线陡）
  \draw[thick, MidBlue] (0.0,1.8) -- (2.0,0.0);
  \draw[thick, MidBlue] (0.0,3.0) -- (3.33,0.0);
  \draw[thick, MidBlue] (0.0,4.5) -- (5.0,0.0);
  % 预算线：斜率 -p1/p2 = -0.5 更平，与最高可达无差异曲线交于 x1 最大的端点
  \draw[thick, dashed] (0.0,2.5) -- (5.0,0.0);
  \fill (5.0,0.0) circle (2pt) node[below right=2pt] {\small 角点解};
  \node[align=center,text width=5.4cm,font=\scriptsize] at (2.5,-1.35) {无差异曲线为平行直线族（斜率 $-a/b=-0.9$），预算线（虚线）更平（$-p_1/p_2=-0.5$）。因 $a/b>p_1/p_2$，商品 1 的边际价值相对更高，最优解不出现在切点，而落在预算线的端点 $(5,0)$ 上：全部收入购买商品 1。};
\end{tikzpicture}
\end{center}
\figexplan{无差异曲线是一族平行直线，代表边际替代率恒定；预算线（虚线）斜率与之不同，因此最优解不是切点而落在预算线的端点（角点解，黑点）。这正是完全替代偏好的标志：不存在内点切解，需求对相对价格极度敏感。}

\notehead{适用情形}
近似同质商品（两个品牌的同规格产品、不同币种的现金）；
也是\textbf{跨期消费}中“今天消费 vs 明天消费”极简模型的偏好形式。
不适用：任何需要平滑替代、内点解的分析。

\subsection{完全互补（固定比例效用）}

\notehead{形式}
$u(x_1,x_2)=\min\{ax_1,bx_2\}$（$a,b>0$）。

\notehead{性质与推导}
无差异曲线为 L 形，拐点在射线 $ax_1=bx_2$ 上；$\sigma=0$。
最优必在拐点（拐点之外多买的商品不增效用却耗预算）：
\[
ax_1=bx_2\ \Longrightarrow\ x_2=\frac{a}{b}x_1;\qquad
p_1x_1+p_2\frac{a}{b}x_1=m
\]
\[
\boxed{\ x_1^*=\frac{mb}{bp_1+ap_2},\qquad x_2^*=\frac{ma}{bp_1+ap_2}\ }
\]
间接效用 $v(p,m)=\dfrac{mab}{bp_1+ap_2}$；支出函数 $e(p,u)=\dfrac{u(bp_1+ap_2)}{ab}$。
需求只依赖“组合价格” $(bp_1+ap_2)/(ab)$：单一商品涨价与组合整体涨价效果相同。

\begin{center}
\begin{tikzpicture}[scale=0.9]
  \draw[axisline] (0,0) -- (5.4,0) node[below left, axislabel] {$x_1$};
  \draw[axisline] (0,0) -- (0,4.4) node[below left, axislabel] {$x_2$};
  \draw[dashed, MidBlue] (0,0) -- (2.3,4.6) node[right] {$ax_1=bx_2$};
  \draw[thick, MidBlue] (0,3.0) -- (1.5,3.0) -- (1.5,0);
  \draw[thick, MidBlue] (0,2.0) -- (1.0,2.0) -- (1.0,0);
  \draw[thick, MidBlue] (0,4.0) -- (2.0,4.0) -- (2.0,0);
  \draw[thick, dashed, MidBlue] (0.3,3.9) -- (4.7,0.8);
  \node[nodeblue] at (1.5,3.0) {};
  \node[align=left] at (3.7,2.0) {L 形无差异曲线\\ 最优在拐点（黑点）\\ 无替代余地};
\end{tikzpicture}
\end{center}
\figexplan{无差异曲线为 L 形（直角），其拐点全部落在过原点的射线 $ax_1=bx_2$ 上——该射线即固定比例的搭配轨迹。最优解位于预算线（虚线）与该射线的交点所对应的拐点（黑点）：两种商品必须按比例搭配，没有任何替代余地。}

\notehead{适用情形}
严格固定比例：左脚鞋与右脚鞋、发动机与车身、打印作业中的纸与墨；
生产理论中的固定比例技术是其对偶。也是\textbf{悲观预期}（maximin）偏好的几何原型。

\subsection{乘积型效用（Cobb--Douglas 型）}

\notehead{形式}
$u(x_1,x_2)=x_1^{\alpha}x_2^{\beta}$（$\alpha,\beta>0$）。

\notehead{性质与推导}
取严格递增变换 $\ln$，等价于对数可加形式 $\alpha\ln x_1+\beta\ln x_2$（第三章 §4.5：表示不唯一）。
$\sigma=1$，位似，支出份额恒定。一阶条件
\[
\mathrm{MRS}=\frac{\alpha x_2}{\beta x_1}=\frac{p_1}{p_2}
\ \Longrightarrow\ \beta p_1x_1=\alpha p_2x_2,
\]
代入预算：$\alpha(p_1x_1+p_2x_2)=\alpha m$，而 $p_2x_2=(\beta/\alpha)p_1x_1$，
\[
\boxed{\ x_1^*=\frac{\alpha}{\alpha+\beta}\frac{m}{p_1},\qquad x_2^*=\frac{\beta}{\alpha+\beta}\frac{m}{p_2}\ }
\]
\textbf{每种商品的支出份额恒为权重占比}（$\alpha/(\alpha+\beta)$），与价格、收入无关——
这是该形式被大量使用的全部理由，也是被批评的全部理由。
间接效用（记 $\bar\alpha=\alpha/(\alpha+\beta)$ 等）：
$v(p,m)=\bar\alpha^{\alpha}\bar\beta^{\beta}\,m^{\alpha+\beta}p_1^{-\alpha}p_2^{-\beta}$。

\begin{center}
\begin{tikzpicture}[scale=0.9]
  \draw[axisline] (0,0) -- (5.4,0) node[below left, axislabel] {$x_1$};
  \draw[axisline] (0,0) -- (0,4.4) node[below left, axislabel] {$x_2$};
  \draw[thick, MidBlue, domain=0.55:5.0, smooth, variable=\x] plot ({\x},{1.7/\x});
  \draw[thick, MidBlue, domain=0.9:5.0, smooth, variable=\x] plot ({\x},{3.0/\x});
  \draw[thick, MidBlue, domain=1.35:5.0, smooth, variable=\x] plot ({\x},{4.3/\x});
  \draw[thick, dashed, MidBlue] (0.4,3.7) -- (4.9,0.9);
  \node[nodeblue] at (2.11,1.42) {};
  \node[align=left] at (3.8,3.4) {双曲线族\\ 份额恒定 $\Rightarrow$\\ 收入扩张路径为直线};
  \draw[dashed, MidBlue] (0,0) -- (3.4,2.29);
\end{tikzpicture}
\end{center}
\figexplan{无差异曲线为双曲线族。由于每种商品的支出份额恒定，收入扩张路径（自原点出发的虚线射线）是一条直线；另一条虚线为预算线，其与无差异曲线的切点即最优解。份额恒定也意味着恩格尔曲线为过原点的直线。}

\notehead{适用情形}
宏观基准模型（消费—闲暇、跨期对数效用 $\ln c$ 是其特例）、
需要\textbf{闭式解}的任何教学与基准分析；
经验上适用于支出份额近似稳定的加总品类。
不适用：份额随收入系统变化（必需品/奢侈品）的情形——那要 §\ref{subsec:sg} 或非位似形式。

\subsection{常替代弹性效用（CES）}

\notehead{形式}
$u(x_1,x_2)=\bigl(\alpha x_1^{\rho}+\beta x_2^{\rho}\bigr)^{1/\rho}$（$\rho\le1,\ \rho\ne0$）。

\notehead{性质与推导}
一次齐次（位似）。边际替代率
\[
\mathrm{MRS}=\frac{\alpha x_1^{\rho-1}}{\beta x_2^{\rho-1}}=\frac{\alpha}{\beta}\Bigl(\frac{x_1}{x_2}\Bigr)^{\rho-1},
\qquad
\sigma=\frac{\mathrm{d}\ln(x_2/x_1)}{\mathrm{d}\ln\mathrm{MRS}}=\frac{1}{1-\rho}.
\]
切线条件 $\mathrm{MRS}=p_1/p_2$ 给出
\[
\frac{x_1}{x_2}=\Bigl(\frac{\alpha}{\beta}\frac{p_2}{p_1}\Bigr)^{\sigma},
\]
代回预算约束，得需求与价格指数：
\[
\boxed{\ x_i^*=\frac{\alpha_i^{\sigma}p_i^{-\sigma}}{\alpha^{\sigma}p_1^{1-\sigma}+\beta^{\sigma}p_2^{1-\sigma}}\,m,
\qquad
P\equiv\bigl(\alpha^{\sigma}p_1^{1-\sigma}+\beta^{\sigma}p_2^{1-\sigma}\bigr)^{\frac{1}{1-\sigma}},
\qquad v(p,m)=\frac{m}{P}\ }
\]
（$\alpha_1=\alpha,\alpha_2=\beta$；检验 $\sigma=1$：$x_i^*=\alpha_i m/p_i$，回到乘积型。）

\textbf{三个极限}（谱系图的严格内容）：
\begin{itemize}
  \item $\rho\to1$（$\sigma\to\infty$）：$u\to\alpha x_1+\beta x_2$，完全替代；
  \item $\rho\to0$（$\sigma\to1$）：对 $\ln u=\frac{1}{\rho}\ln(\alpha x_1^{\rho}+\beta x_2^{\rho})$
        用洛必达，$\lim_{\rho\to0}\ln u=\frac{\alpha\ln x_1+\beta\ln x_2}{\alpha+\beta}$，乘积型；
  \item $\rho\to-\infty$（$\sigma\to0$）：括号内被较小的 $x_i$ 主导，$u\to\min$ 型，完全互补。
\end{itemize}

\begin{center}
\begin{tikzpicture}[scale=0.9]
  \draw[axisline] (0,0) -- (5.6,0) node[below left, axislabel] {$x_1$};
  \draw[axisline] (0,0) -- (0,4.4) node[below left, axislabel] {$x_2$};
  % sigma large (near linear)
  \draw[thick, MidBlue] (0.7,3.5) -- (4.9,0.75);
  % CD
  \draw[thick, MidBlue, domain=0.75:5.0, smooth, variable=\x] plot ({\x},{2.6/\x});
  % low sigma (near L): smooth curve close to L
  \draw[thick, MidBlue] (1.05,3.8) .. controls (1.1,1.5) and (1.45,1.15) .. (4.6,0.85);
  \node[align=left] at (4.3,3.7) {曲率 $=\sigma$ 的反面：\\ 近直线（$\sigma$ 大）\\ 双曲（$\sigma=1$）\\ 近 L 形（$\sigma$ 小）};
\end{tikzpicture}
\end{center}
\figexplan{三条曲线显示 CES 族随替代弹性 $\sigma$ 变化的形态：$\sigma$ 大时接近直线（趋向完全替代），$\sigma=1$ 时为双曲线（乘积型），$\sigma$ 小时逼近 L 形（趋向完全互补）。曲率正是替代弹性的反面：越弯曲，两商品越难相互替代。}

\notehead{适用情形}
需要\textbf{参数化替代弹性}的一切场合：国际贸易（国家间产品替代）、
宏观加总（资本—劳动）、垄断竞争（品种间替代）、能源模型。
其垄断竞争特例 $\sigma>1$ 且品种数内生，是贸易与空间经济学的标准件。

\subsection{拟线性效用}

\notehead{形式}
$u(x_1,x_2)=v(x_1)+x_2$，$x_2$ 为计价物。

\notehead{性质与推导}
第三章 §4.7 已给出偏好侧刻画（无财富效应）。需求侧：
$x_2$ 为计价物（$p_2=1$），内点一阶条件
\[
v'(x_1^*)=p_1\ \Longrightarrow\ x_1^*=(v')^{-1}(p_1),
\]
\textbf{与收入 $m$ 无关}（直到 $m$ 不足以购买 $(v')^{-1}(p_1)$ 时出现角点）。
例 $v(x_1)=\sqrt{x_1}$：$\frac{1}{2\sqrt{x_1}}=p_1\Rightarrow x_1^*=\frac{1}{4p_1^2}$，
需求弹性恒为 $-2$。收入扩张路径为竖直线；消费者剩余
$\int_{p_1}^{\infty}x_1^*(q)\,\mathrm{d}q$ 是福利变化的\textbf{精确}度量（无收入效应扭曲）。

\begin{center}
\begin{tikzpicture}[scale=0.9]
  \draw[axisline] (0,0) -- (5.4,0) node[below left, axislabel] {$x_1$};
  \draw[axisline] (0,0) -- (0,4.4) node[below left, axislabel] {$x_2$（计价物）};
  \draw[thick, MidBlue, domain=0.02:5.0, smooth, variable=\x] plot ({\x},{3.5-sqrt(\x)});
  \draw[thick, MidBlue, domain=0.02:5.0, smooth, variable=\x] plot ({\x},{2.6-sqrt(\x)});
  \draw[thick, MidBlue, domain=0.02:5.0, smooth, variable=\x] plot ({\x},{1.7-sqrt(\x)});
  \draw[dashed, thick] (2.1,0) -- (2.1,4.2) node[above] {收入扩张路径};
  \node[align=left] at (4.15,3.6) {无差异曲线纵向平移\\ 最优 $x_1$ 不随收入变};
\end{tikzpicture}
\end{center}
\figexplan{无差异曲线互为纵向平移，即在同一 $x_1$ 处斜率恒为 $v'(x_1)$。收入扩张路径是一条竖直线：一阶条件 $v'(x_1)=p_1$ 不含收入，故最优 $x_1^*$ 与收入无关，收入变化全部由计价物吸收。这也是消费者剩余成为福利变化精确度量的原因。}

\notehead{适用情形}
部分均衡分析（单个市场 + 其余一切合成计价物）、拍卖与机制设计、
产业组织、公共经济学的福利计算。
不适用：所研究商品支出占收入比重大（住房、食品之于低收入群体）——
此时收入效应是一阶的，拟线性假设会抹掉它。

\subsection{基本需求型（Stone--Geary / 线性支出系统）}\label{subsec:sg}

\notehead{形式与推导}
\[
u(x_1,x_2)=(x_1-\gamma_1)^{\alpha}(x_2-\gamma_2)^{\beta},\qquad x_i>\gamma_i\ge0.
\]
$\gamma_i$ 为\textbf{基本需求}（生存/必需量）：先购买 $\gamma_i$，再对
\textbf{剩余收入}（supernumerary income）$\tilde m=m-p_1\gamma_1-p_2\gamma_2$
按乘积型分配。令 $z_i=x_i-\gamma_i$，问题化为 $\max z_1^{\alpha}z_2^{\beta}\ \st p_1z_1+p_2z_2=\tilde m$，
由 §4.1.4（乘积型结果）：
\[
\boxed{\ x_i^*=\gamma_i+\frac{\alpha_i}{\alpha+\beta}\,\frac{m-p_1\gamma_1-p_2\gamma_2}{p_i}\ }
\]
需求对 $(p,m)$ \textbf{线性}（线性支出系统，LES），
恩格尔曲线是直线但\textbf{不过原点}——必需品份额随收入递减、奢侈品份额递增，
恰好补上乘积型“份额恒定”的缺口。

\begin{center}
\begin{tikzpicture}[scale=0.9]
  \draw[axisline] (0,0) -- (5.6,0) node[below left, axislabel] {$x_1$};
  \draw[axisline] (0,0) -- (0,4.4) node[below left, axislabel] {$x_2$};
  \draw[dashed, MidBlue] (1.3,0) -- (1.3,4.2); \draw[dashed, MidBlue] (0,1.0) -- (5.2,1.0);
  \node[below] at (1.3,0) {$\gamma_1$}; \node[left] at (0,1.0) {$\gamma_2$};
  \draw[thick, MidBlue, domain=1.75:5.2, smooth, variable=\x] plot ({\x},{1.0+1.5/(\x-1.3)});
  \draw[thick, MidBlue, domain=2.05:5.2, smooth, variable=\x] plot ({\x},{1.0+2.6/(\x-1.3)});
  \node[nodeblue, label={above right:{渐近原点 $(\gamma_1,\gamma_2)$}}] at (1.3,1.0) {};
  \node[align=left] at (4.35,3.4) {无差异曲线以\\ $(\gamma_1,\gamma_2)$ 为\\ “新原点”};
\end{tikzpicture}
\end{center}
\figexplan{两条虚线标出基本需求量 $(\gamma_1,\gamma_2)$，它构成无差异曲线族的“新原点”：所有曲线都以该点为渐近中心向外弯曲。因此收入扩张路径是直线但不过原点，恩格尔曲线线性而截距非零——必需品份额随收入递减、奢侈品份额递增，补上了乘积型份额恒定的缺口。}

\notehead{适用情形}
发展经济学（食品必需量）、需求系统估计（LES 及其推广几乎理想需求系统）、
跨期模型中的“习惯存量” $\gamma$（消费习惯如同基本需求，§\ref{subsec:habit}）。

\subsection{非凸偏好（平方和效用）}

\notehead{形式}
$u(x_1,x_2)=x_1^{2}+x_2^{2}$。

\notehead{性质}
无差异曲线是以原点为圆心的四分之一圆弧，上轮廓集\textbf{非凸}
（圆周外的区域不含其边界点的均值）：
$(1,0)$ 与 $(0,1)$ 效用均为 $1$，但均值 $(0.5,0.5)$ 的效用为 $0.5<1$——
\textbf{平均消费劣于极端消费}。最优解恒为角点：
预算线上效用沿 $x_1^2+(m-p_1x_1)^2/p_2^2$ 是 $x_1$ 的凸函数，最大值在端点。
$$x^*=\arg\max\bigl\{(m/p_1)^2,\ (m/p_2)^2\bigr\}\ \text{二选一}.$$

\begin{center}
\begin{tikzpicture}[scale=0.9]
  \draw[axisline] (0,0) -- (5.2,0) node[below left, axislabel] {$x_1$};
  \draw[axisline] (0,0) -- (0,4.4) node[below left, axislabel] {$x_2$};
  \draw[thick, MidBlue] (3.6,0) arc (0:90:3.6);
  \draw[thick, MidBlue] (2.4,0) arc (0:90:2.4);
  \draw[thick, dashed, MidBlue] (4.4,0) -- (0,2.9);
  \node[nodeblue, label={above right:{角点最优}}] at (4.4,0) {};
  \node[align=left] at (3.6,3.4) {圆弧族向外弯\\ 均值点劣于端点\\ $\Rightarrow$ 专业化};
\end{tikzpicture}
\end{center}
\figexplan{无差异曲线是向外弯的圆弧族，代表非凸偏好：上轮廓集是圆弧以外的区域，因此平均消费束劣于两个极端端点，最优恒为角点解（黑点）。这一形态刻画专业化与非此即彼的选择。}

\notehead{适用情形}
\textbf{专业化}与\textbf{非此即彼}的选择：职业选择（全时从商或从政）、
成瘾品、网络效应下的标准之争。是凸偏好框架的基准反例。

\begin{figure}[H]
\centering
\begin{tikzpicture}[font=\scriptsize]
  \draw[axisline] (0,0) -- (4.6,0) node[below right] {$x_1$};
  \draw[axisline] (0,0) -- (0,3.6) node[above left] {$x_2$};
  \draw[very thick] (0,2.0) -- (1.4,2.0) -- (1.4,3.0) -- (2.8,3.0) -- (2.8,4.0) -- (4.4,4.0);
  \node[font=\scriptsize] at (3.5,1.35) {(a) 完全可加 $u=u_1(x)+u_2(y)$：阶梯形};
  \begin{scope}[xshift=6.4cm]
    \draw[axisline] (0,0) -- (4.6,0) node[below right] {$x_1$};
    \draw[axisline] (0,0) -- (0,3.6) node[above left] {$x_2$};
    \draw[very thick,domain=0.5:4.3,smooth,variable=\t] plot ({\t},{3.4/\t});
    \draw[very thick,domain=0.9:4.3,smooth,variable=\t] plot ({\t},{5.6/\t});
    \node[font=\scriptsize] at (3.4,1.05) {(b) 不可加（Cobb--Douglas）：内弯双曲线};
  \end{scope}
  \node[align=center,font=\scriptsize] at (3.2,-1.15) {(a) 的无差异曲线是轴对齐的阶梯：两个坐标可分别“换算”，边际替代率在阶梯上无定义；\\ (b) 的无差异曲线光滑内弯， MRS 随 $x_2$ 递减。};
\end{tikzpicture}
\caption{可加与不可加：无差异曲线的形状差异}
\end{figure}
\figexplan{左面板是完全可加效用 $u=u_1(x_1)+u_2(x_2)$：沿无差异曲线，$x_1$ 增加多少可以按固定“汇率”换回 $x_2$，故曲线呈轴对齐的阶梯形，直角处边际替代率无定义。右面板是不可加的乘积型效用：无差异曲线是光滑内弯的双曲线，MRS 随 $x_2$ 递减。正是这一形状差异，使“坐标可分离”成为可检验的行为含义。}

\subsection{加法可分效用}

\notehead{形式}
$u(x)=\sum_{i=1}^{n}u_i(x_i)$。

\notehead{性质}
第三章 §4.6 已给出偏好侧充要条件（坐标独立性 + 矩形条件）。
需求侧含义：交叉偏导 $\partial^2u/\partial x_i\partial x_j=0$，
一种商品的边际效用不依赖其他商品消费量；
在预算约束下，一阶条件 $u_i'(x_i)=\lambda p_i$ 使各商品需求\textbf{仅通过乘子 $\lambda$ 联动}
（收入效应渠道），无直接交叉替代。
乘积型的对数形式 $\sum\alpha_i\ln x_i$、CES 的单调变换、拟线性都是其特例。

\notehead{适用情形}
跨期效用 $\sum_t\delta^t u(c_t)$ 的静态影子；多商品需求系统；
需要“边际效用互不干扰”结构来识别参数的经验模型。

\subsection{柔性函数形式（需求估计用）}\label{subsec:flexible}

\begin{center}
\footnotesize
\setlength{\tabcolsep}{4pt}
\begin{tabular}{llll}
\toprule
名称 & 形式 & 性质 & 适用 \\
\midrule
超越对数 & $\ln u=\alpha_0+\sum_i\alpha_i\ln x_i+\frac12\sum_{ij}\beta_{ij}\ln x_i\ln x_j$ & 二阶近似、弹性可变 & 需求系统 \\
广义固定比例 & $u=\sum_{ij}b_{ij}\sqrt{x_ix_j}$（$b_{ij}=b_{ji}$） & 一次齐次、柔性 & 成本对偶 \\
极式 & 间接效用 $F(p)\cdot v(x)+G_i(p)$ 型 & 恩格尔曲线线性 & 可加总 \\
分支树型 & 嵌套 CES（分支内 $\rho_k$、分支间 $\sigma_k$） & 分组替代、层级 & 多品类需求 \\
\bottomrule
\end{tabular}
\end{center}

\notehead{要点}
柔性形式不追求“有经济学名字的偏好”，而是追求\textbf{对任意弹性结构的二阶近似能力}：
参数多到足以让弹性随点变化。代价是参数多、需要大数据、
必须事后检验正则性（拟凹、单调）是否在样本域内成立。
分支树型是“CES 的 CES”：例如先在外层选择食品/服装/服务，
再在外层之下做品类内替代——正是离散选择中嵌套结构（§\ref{subsec:rum}）的连续对偶。

% ============================================================
\section{需求系统与更多静态形式（补全）}
% ============================================================

\subsection{PIGLOG 与几乎理想需求系统（AIDS）}

\notehead{为什么需要它}
CES 族（§4.1.5）恩格尔曲线迫使份额对收入不敏感；
基本需求型（§4.1.7）只允许份额随收入线性衰减。
经验家计数据反复显示：\textbf{份额—收入关系既非平线也非直线}。
AIDS 的做法是：不先写效用，先写\textbf{支出函数}，
让它对价格与效用都足够灵活，再反推需求。

\notehead{形式（PIGLOG 支出函数）}
\[
\ln e(p,u)=\underbrace{\alpha_0+\sum_k\alpha_k\ln p_k+\frac12\sum_k\sum_j\gamma_{kj}^{*}\ln p_k\ln p_j}_{a(p)}
+\ u\,\underbrace{\beta_0\prod_k p_k^{\beta_k}}_{b(p)}.
\]
效用通过谢泼德引理进入需求：对 $\ln e$ 关于 $\ln p_i$ 求导，
左边 $=\frac{p_i}{e}\frac{\partial e}{\partial p_i}=\frac{p_ix_i}{e}=w_i$（支出份额），故
\[
\boxed{\ w_i=\alpha_i+\sum_j\gamma_{ij}\ln p_j+\beta_i\ln\bigl(m/P\bigr)\ }
\]
其中 $\gamma_{ij}=(\gamma_{ij}^{*}+\gamma_{ji}^{*})/2$，
$\ln P$ 为 $a(p)$ 给出的价格指数（实证中常用其线性近似）。

\notehead{理论约束（效用最大化可积性）}
\begin{center}
\small
\begin{tabular}{lll}
\toprule
约束 & 参数表达 & 含义 \\
\midrule
加总（adding-up） & $\sum_i\alpha_i=1,\ \sum_i\beta_i=0,\ \sum_i\gamma_{ij}=0$ & 份额和为 1 \\
齐次（homogeneity） & $\sum_j\gamma_{ij}=0\ (\forall i)$ & 无货币幻觉 \\
斯勒茨基对称 & $\gamma_{ij}=\gamma_{ji}$ & 可积回效用 \\
\bottomrule
\end{tabular}
\end{center}

\notehead{性质：为什么选它}
(i) 它是任意需求系统的\textbf{一阶柔性近似}；
(ii) 份额方程对参数\textbf{线性}——可加个体固定效应、可用面板工具变量，
这是它在应用微观中统治地位的技术基础；
(iii) 三条理论约束都是\textbf{参数的线性等式}，
检验“需求是否服从效用最大化”就是三个 Wald 检验；
(iv) 恩格尔曲线：$w_i$ 对 $\ln m$ 线性 $\Rightarrow$ 必需品（$\beta_i<0$）
与奢侈品（$\beta_i>0$）份额随收入单调渐变——比 CES 灵活，但仍是直线。

\notehead{适用与案例}
食品、能源等多品类家计需求；福利核算（$e(p,u)$ 已知 $\Rightarrow$
真实价格指数、补偿变化都有公式）。案例：估计中国城镇家庭八大类消费，
食品份额对 $\ln m$ 的斜率显著为负、文化娱乐显著为正，
且齐次性在改革期样本上被拒绝（价格体系剧变期的货币幻觉/调整滞后）。

\subsection{QUAIDS：二次恩格尔曲线}

\notehead{形式}
在 AIDS 份额方程上加二次项（$b(p)=\prod p_j^{\beta_j}$ 为协价格指数的修正）：
\[
w_i=\alpha_i+\sum_j\gamma_{ij}\ln p_j+\beta_i\ln\Bigl(\frac{m}{a(p)}\Bigr)
+\frac{\lambda_i}{b(p)}\Bigl[\ln\Bigl(\frac{m}{a(p)}\Bigr)\Bigr]^{2},
\qquad \sum_i\lambda_i=0.
\]
约束加总/齐次/对称同 AIDS，仅多 $\sum_i\lambda_i=0$。

\notehead{性质：为什么选它}
AIDS 的恩格尔曲线是直线，意味着“必需品永远是必需品”；
QUAIDS 的二次项允许\textbf{曲线弯折}：同一种商品可以对穷人是奢侈品、
对富人是必需品（肉制品、在外就餐的标准形态）。
理论意义：可积性仍成立（存在对应的间接效用），
是“最弯的仍理论自洽”的需求系统。
近年应用（联网核实）：冈比亚家庭食品价格福利分析、
美国食品需求弹性的理论一致元分析——QUAIDS 都是首选，
理由正是其非线性恩格尔结构与约束可检验性。

\begin{center}
\begin{tikzpicture}[scale=0.95]
  \draw[axisline] (0,0) -- (6.6,0) node[below left, axislabel] {$\ln m$};
  \draw[axisline] (0,0) -- (0,4.2) node[below left, axislabel] {份额 $w_i$};
  \draw[thick, MidBlue] (0.3,3.4) -- (5.6,1.3) node[above right, align=left] {必需品（AIDS：直线）};
  \draw[thick, MidBlue] (0.3,0.9) .. controls (2.2,1.2) and (3.6,2.6) .. (5.4,3.5)
    node[above right, align=left] {QUAIDS：先平后陡};
  \draw[thick, MidBlue] (0.3,2.9) .. controls (2.4,2.5) and (3.9,1.5) .. (5.2,0.5)
    node[below right, align=left] {QUAIDS：奢侈品转必需品};
  \node[align=center] at (3.2,-0.9) {二次项使份额—收入关系可弯折：同一商品在不同收入档角色可变};
\end{tikzpicture}
\end{center}
\figexplan{横轴为收入对数 $\ln m$，纵轴为支出份额 $w_i$。AIDS 给出直线（份额与 $\ln m$ 线性）；QUAIDS 引入二次项后份额—收入关系可以弯折：同一商品可在低收入档表现为必需品、在高收入档表现为奢侈品（或相反）。这正是 QUAIDS 能拟合非线性恩格尔结构的原因。}

\subsection{非位似 CES（隐式加法可分）}

\notehead{形式与核心性质}
效用 $U$ 由隐式方程定义：
\[
\sum_i \Omega_i\,\bigl(c_i\,U^{\varepsilon_i}\bigr)^{\frac{\sigma-1}{\sigma}}=1,
\qquad \Omega_i>0,\ \sigma>0,
\]
$\varepsilon_i$ 为部门 $i$ 的收入弹性参数。
直接对隐函数求导可得相对恩格尔曲线的斜率：
\[
\boxed{\ \frac{\partial\ln(c_i/c_j)}{\partial\ln C}=\varepsilon_i-\varepsilon_j\ }
\]
即\textbf{相对支出对对数收入是直线，且斜率恒定}。
对比：位似 CES（斜率 0）、基本需求型（斜率随收入快速衰减至 0）。
跨国与跨期数据显示相对份额—相对收入的关系恰是稳定直线
（农业下降、服务业上升的全谱系）——这是该形式被设计出来要拟合的靶心。

\notehead{为什么选它}
研究结构转型、消费升级、创新方向（需求拉动）时，
需要长期、大幅收入变动下\textbf{不衰减}的收入效应；
该形式是目前唯一同时做到“任意部门数、收入与价格效应并存、
可估计”的偏好族。代价：无显式需求函数，数值求解。

\subsection{家庭生产与时间配置}

\notehead{形式}
\[
U=U(Z_1,\dots,Z_K),\qquad Z_k=f_k(x_k,t_k),
\]
效用不直接来自市场商品 $x_k$，而来自\textbf{家庭内生产的商品} $Z_k$
（营养、清洁、亲子陪伴），其投入是市场品 $x_k$ 与时间 $t_k$。
预算：$\sum_k p_kx_k=wT_w+V$（$T_w$ 工时、$V$ 非劳动收入）；
时间：$\sum_k t_k+T_w=T$。
合并得\textbf{完全收入}约束
\[
\sum_k\bigl(p_kx_k+wt_k\bigr)=wT+V\equiv S,
\]
商品 $Z_k$ 的\textbf{影子价格} $\pi_k$ 含时间成本：
$\pi_k=\min_{x,t}\{px+wt:\ f_k(x,t)\ge1\}$。

\notehead{为什么选它：能反映什么}
工资上升同时抬高市场品与时间的价格（后者以 $w$ 计），
故\textbf{时间密集商品}（照料、休闲活动）相对变贵——
解释高收入社会的时间荒、家务外包、生育率下降
（子女是时间最密集的 $Z$）。
这也是 §4.9 生育模型把养育成本写为 $v\,w$（工资的比例）的理论依据：
养育成本中时间投入随工资升值。

\subsection{利他、暖光与公共品私人供给}

\notehead{三种写法}
\begin{enumerate}
  \item \textbf{纯利他}：$u_i=u_i(x_i,G)$，$G=\sum_j g_j$ 为公共品总量
        ——只在意结果；
  \item \textbf{不纯利他（暖光）}：$u_i=u_i(x_i,G,g_i)$，自己的捐献 $g_i$ 单独进入效用
        ——在意“我捐了”本身；
  \item \textbf{消费—捐赠线性型}：$u_i=x_i+\gamma_i\ln G$（常用可解形式）。
\end{enumerate}

\notehead{关键结论（推导）}
纯利他 + 内点解下，把他人捐赠 $\sum_{j\ne i}g_j$ 视为收入调整，
个体最优 $G$ 只依赖总收入分布\textbf{之外}的总量：
政府以税收融资的公共供给增量 $\Delta G_{gov}$
被私人捐赠等额挤出（dollar-for-dollar crowding-out）。
暖光项打破完全挤出：自己的 $g_i$ 有独立边际效用，
挤出率小于 1——实验与捐赠数据支持不纯利他。

\notehead{适用}
慈善、众筹、家庭内部转移（§4.7.9 中国赡养转移的模型化）、ESG 消费。

\subsection{议价：合作剩余的分割}

\notehead{形式}
可行效用对的凸集 $S$，破裂点 $d=(d_1,d_2)$。乘积型目标
\[
\max_{(u_1,u_2)\in S,\ u\ge d}\ (u_1-d_1)(u_2-d_2)\ \ \text{（几何平均，等价于取对数后加法可分）}.
\]
\textbf{推导（可转移效用的特例）}：设前沿 $u_1+u_2=\bar S$，代入：
\[
\max_{u_1}\ (u_1-d_1)(\bar S-u_1-d_2)
\ \Longrightarrow\
u_1^{*}-d_1=u_2^{*}-d_2=\frac{\bar S-d_1-d_2}{2}:
\]
\textbf{剩余均分}。一般前沿 $u_2=g(u_1)$ 下，一阶条件
$-g'(u_1^{*})=\dfrac{u_2^{*}-d_2}{u_1^{*}-d_1}$：
前沿斜率 = 无差异双曲线的斜率（切点条件）。

\notehead{为什么选它}
把“谈判力”参数化为破裂点 $d$ 与权重
（加权乘积 $(u_1-d_1)^{\eta}(u_2-d_2)^{1-\eta}$），
闭式、可嵌入家庭（离婚威胁点）、劳资（罢工威胁点）、
国际贸易（协定外选项）模型。
家庭经济学中，集体模型 + 议价解是分析夫妻间分配
（生育决策、退休协调）的标准工具——与 §4.9 的家庭决策直接衔接。

% ============================================================
\section{跨期与风险效用函数}
% ============================================================

\begin{figure}[H]
\centering
\begin{tikzpicture}[font=\scriptsize]
  \draw[axisline] (0,0) -- (6.0,0) node[below right] {今天消费 $c_1$};
  \draw[axisline] (0,0) -- (0,4.2) node[above left] {明天消费 $c_2$};
  % 跨期预算线
  \draw[thick] (0,3.2) -- (4.9,0) node[midway,above left] {跨期预算线 $c_2+\frac{c_1}{1+r}=w$};
  % 禀赋点与借贷
  \draw[dashed,thin] (2.2,1.05) -- (2.2,0) node[below=2pt] {禀赋点};
  \node[font=\scriptsize] at (2.2,0.72) {借贷点};
  % 无差异曲线
  \draw[very thick,domain=0.35:3.6,smooth,variable=\t] plot ({\t},{2.6-0.45*\t*\t});
  \draw[very thick,domain=1.9:4.6,smooth,variable=\t] plot ({\t},{3.5-0.16*\t*\t});
  \fill (2.2,1.55) circle (1.5pt) node[above right=2pt] {$c^{*}$};
  \node[align=center,font=\scriptsize] at (4.6,3.35) {无差异曲线内弯（凸偏好）\\最优点为切点，取决于利率 $r$};
\end{tikzpicture}
\caption{跨期选择：两期无差异曲线、跨期预算线与借贷点}
\end{figure}
\figexplan{横轴为今天消费、纵轴为明天消费。跨期预算线斜率为 $-(1+r)$，即相对价格；虚线标出禀赋点，它把平面分为“借贷区”（今日消费大于禀赋）与“储蓄区”。无差异曲线内弯表示偏好具凸性，最优点 $c^{*}$ 是无差异曲线与跨期预算线的切点，满足 $u'(c^{*}_1)=\beta u'(c^{*}_2)(1+r)$。图形要点：利率改变的是斜率，从而改变储蓄与借贷的方向。}

\subsection{指数贴现效用}\label{subsec:expdisc}

\notehead{形式}
$U=\sum_{t=0}^{T}\delta^{t}u(c_t)$（$\delta\in(0,1)$）。

\notehead{性质与推导}
时间可加 + 贴现率恒定：$t$ 与 $t+1$ 之间的边际替代率
$\mathrm{MRS}_{t,t+1}=\delta\,u'(c_{t+1})/u'(c_t)$ 只依赖两期消费，
与决策时点无关——这是\textbf{动态一致性}的来源：
今天为明天做的计划，明天不会推翻。
取 $u(c)=c^{1-\gamma}/(1-\gamma)$（$\gamma>0$），跨期欧拉方程
\[
u'(c_t)=\delta(1+r_{t+1})\,u'(c_{t+1})
\ \Longrightarrow\
c_t^{-\gamma}=\delta(1+r_{t+1})\,c_{t+1}^{-\gamma}
\ \Longrightarrow\
\boxed{\ \frac{c_{t+1}}{c_t}=\bigl[\delta(1+r_{t+1})\bigr]^{1/\gamma}\ }
\]
消费增长率由\textbf{跨期替代弹性} $\mathrm{EIS}=1/\gamma$ 调节：
$\gamma$ 越大，消费路径越平（不愿在不同时期间挪动消费）。
$\gamma=1$ 时 $u=\ln c$，储蓄率对利率不变（收入效应与替代效应相消）。

\notehead{适用情形}
一切动态基准：增长模型、生命周期储蓄、消费平滑。
局限：EIS 与相对风险厌恶被同一个 $\gamma$ 绑定（见 §\ref{subsec:ez} 的解绑）。

\subsection{习惯形成}\label{subsec:habit}

\notehead{内部习惯}
\[
u(c_t,h_t)=\frac{(c_t-\theta h_t)^{1-\gamma}}{1-\gamma},\qquad
h_t=\sum_{s=1}^{\infty}\phi^{s}c_{t-s},
\]
效用依赖\textbf{超出习惯存量的部分}（差值形式）：
今天的消费抬高明天的参照标准，产生“消费棘轮”。
结构上与基本需求型（§\ref{subsec:sg}）同构：$h_t$ 就是内生的 $\gamma$。

\notehead{外部习惯（攀比形式）}
\[
u(c_t,\bar c_t)=\frac{\bigl(c_t/\bar c_t^{\theta}\bigr)^{1-\gamma}}{1-\gamma},
\]
效用依赖自己消费与社会平均消费 $\bar c_t$ 的\textbf{比值}：
个体不把 $\bar c$ 当决策变量，均衡中出现“消费军备竞赛”——
全社会消费都涨，个体效用不涨（位置性外部性），
为\textbf{消费税/累进税}提供规范依据。

\notehead{适用情形}
资产定价（解释股权溢价、消费的过度平滑）、储蓄的持续性、
幸福悖论（收入增长而主观福利停滞）。
差值形式易解（近似对数线性化后保持线性），比值形式政策含义强。

\subsection{递归偏好与 EIS--风险厌恶解绑}\label{subsec:ez}

\notehead{Kreps--Porteus 递归形式}
\[
U_t=W\bigl(c_t,\ \mathcal{I}_t(U_{t+1})\bigr),
\]
当期消费与“未来效用的确定性等价” $\mathcal{I}_t$ 通过聚合器 $W$ 复合。
时间可加指数贴现是其特例：$W(c,z)=u(c)+\delta z$，$\mathcal{I}_t=\mathbb{E}_t$。

\notehead{Epstein--Zin 形式}
\[
U_t=\Bigl[(1-\delta)c_t^{1-\rho}+\delta\Bigl(\mathbb{E}_t\bigl[U_{t+1}^{1-\gamma}\bigr]\Bigr)^{\frac{1-\rho}{1-\gamma}}\Bigr]^{\frac{1}{1-\rho}},
\]
两个参数各司其职：$\mathrm{EIS}=1/\rho$（跨期替代意愿），
$\gamma$ = 相对风险厌恶（对不确定性的态度）。
\textbf{检验归约}：当 $\rho=\gamma$ 时，
\[
U_t^{1-\gamma}=(1-\delta)c_t^{1-\gamma}+\delta\,\mathbb{E}_t\bigl[U_{t+1}^{1-\gamma}\bigr],
\]
即 $U^{1-\gamma}$ 是期望效用型的——退化为时间可加的 CRRA 贴现期望效用。
当 $\gamma>\rho$：偏好\textbf{不确定性早解决}（尽早知道命运）；
$\gamma<\rho$：偏好晚解决。时间可加模型无法表达这种态度（$\gamma=\rho$ 强制中立）。

\notehead{适用情形}
资产定价（长期风险模型）、气候经济学（百年尺度的不确定性+贴现）、
任何需要“风险厌恶很强烈但跨期替代不必同样强烈”的场合——
数据上两者本就不相等，时间可加形式是强加的束缚。

\subsection{准双曲贴现（贝塔-德尔塔形式）}\label{subsec:betadelta}

\notehead{形式}
\[
U_t=u(c_t)+\beta\sum_{s=1}^{\infty}\delta^{s}u(c_{t+s}),\qquad \beta\in(0,1].
\]
“现在”与“一切未来”之间有一个额外折扣 $\beta$：
$t$ 与 $t+1$ 的权衡，从 $t$ 看是 $\beta\delta$，从 $t+1$ 看变成 $\delta$——
\textbf{动态不一致}：计划推翻计划。

\notehead{数值例（计划反转）}
设 $u$ 线性，$\beta=1/2$，$\delta=1$。任务“写报告”：今天做成本 $6$（今天承担），
收益 $10$ 明天到账；明天做成本仍是 $6$（明天承担），收益 $10$ 后天到账。
\begin{itemize}
  \item \textbf{今天的计算}：今天做：$-6+\frac12\cdot10=-1$；
        明天做：$\frac12(-6+10)=2$；不做：$0$。决定：\textbf{明天做}。
  \item \textbf{明天的计算}：明天做：$-6+\frac12\cdot10=-1$；
        后天做：$\frac12\cdot4=2$；不做：$0$。决定：\textbf{再推一天}。
\end{itemize}
每一天都把任务推给下一天：拖延是 $\beta<1$ 的理性结果。
\textit{老练型}个体预见自己的推翻，会寻求\textbf{承诺装置}
（截止期限、自动储蓄扣款）；\textit{天真型}个体反复惊讶。
对照：$\beta=1$（指数贴现）时两天的计算结果一致，无拖延。

\notehead{适用情形}
储蓄不足与养老金设计（自动加入机制）、成瘾、健身与学习投资、
一切“自我控制”问题。经验估计中 $\beta\approx0.5$--$0.8$ 是常见区间。

\subsection{期望效用与风险厌恶的度量}\label{subsec:risk}

\notehead{vNM 期望效用}
$U=\sum_i p_i u(c_i)$，独立性公理是其偏好基础。
风险态度完全由 $u$ 的曲率刻画：
\[
\text{绝对风险厌恶}\ \mathrm{ARA}(w)=-\frac{u''(w)}{u'(w)},\qquad
\text{相对风险厌恶}\ \mathrm{RRA}(w)=-\frac{w\,u''(w)}{u'(w)}.
\]
（推导：对公平赌局 $\tilde\varepsilon$ 收风险金 $\pi$，泰勒展开
$u(w-\pi)\approx\mathbb{E}[u(w+\tilde\varepsilon)]$ 得
$\pi\approx\frac{1}{2}\mathrm{ARA}(w)\cdot\mathrm{Var}(\tilde\varepsilon)$。）

\notehead{三大族系（均为一参数可估）}
\begin{enumerate}
  \item \textbf{CRRA}：$u(w)=\dfrac{w^{1-\gamma}}{1-\gamma}$（$\gamma>0$，$\gamma=1$ 时为 $\ln w$）。
        $u'=w^{-\gamma}$，$u''=-\gamma w^{-\gamma-1}$，
        $\mathrm{ARA}=\gamma/w$（递减），$\mathrm{RRA}=\gamma$（\textbf{常数}）：
        风险金占财富比例不随财富变。宏观金融默认值。
  \item \textbf{CARA}：$u(w)=-e^{-\alpha w}$（$\alpha>0$）。
        $u'=\alpha e^{-\alpha w}$，$u''=-\alpha^2 e^{-\alpha w}$，
        $\mathrm{ARA}=\alpha$（\textbf{常数}）：风险金的绝对额不随财富变。
        正态收益 + CARA $\Rightarrow$ 均值—方差效用
        $\mathbb{E}w-\frac{\alpha}{2}\mathrm{Var}(w)$，组合选择闭式解。
  \item \textbf{二次型}：$u(w)=w-\theta w^2$（$w<1/(2\theta)$）。
        $u'=1-2\theta w$，$\mathrm{ARA}=2\theta/(1-2\theta w)$（\textbf{递增}）：
        越富越怕风险——与直觉相反；且 $w>1/(2\theta)$ 后边际效用为负（餍足点）。
        优点仅是均值—方差结构，现代应用已少，但教学上仍是
        “递增绝对风险厌恶”的标准反例。
\end{enumerate}

\begin{center}
\begin{tikzpicture}[scale=0.86]
  \draw[axisline] (0,0) -- (6.4,0) node[below left, axislabel] {$w$};
  \draw[axisline] (0,-2.6) -- (0,2.6) node[below left, axislabel] {$u(w)$};
  % ln
  \draw[thick, MidBlue, domain=0.12:6.0, smooth, variable=\x] plot ({\x},{ln(\x)});
  \node at (5.6,1.95) {$\ln w$};
  % CRRA gamma=2: 1-1/w
  \draw[thick, MidBlue, domain=0.42:6.0, smooth, variable=\x] plot ({\x},{1-1/\x});
  \node at (3.1,0.95) {$\gamma=2$};
  % CARA: 1-exp(-x)
  \draw[thick, MidBlue, domain=0.02:6.0, smooth, variable=\x] plot ({\x},{1-exp(-\x)});
  \node at (4.7,1.28) {CARA};
  % sqrt
  \draw[thick, MidBlue, domain=0.02:6.0, smooth, variable=\x] plot ({\x},{sqrt(\x)-0.6});
  \node at (5.9,1.45) {$\gamma=1/2$};
  \node[align=left] at (2.2,-2.1) {曲率越大 $\Rightarrow$ 风险厌恶越强；\\ 直线（风险中性）为 $\gamma=0$ 的极限};
\end{tikzpicture}
\end{center}
\figexplan{四条曲线均为递增凹函数，凹性（曲率）越强代表风险厌恶越强：$\ln w$ 与 $\gamma=2$ 属于 CRRA 族，CARA 为常数绝对风险厌恶，$\gamma=1/2$ 曲率最平缓。直线（风险中性）对应 $\gamma=0$ 的极限情形。}

\notehead{HARA 统一表}
三族都是\textbf{双曲绝对风险厌恶}族的成员：
$\mathrm{ARA}(w)=1/(Aw+B)$。
\begin{center}
\small
\begin{tabular}{llll}
\toprule
族 & 参数限制 & ARA 形态 & 经验含义 \\
\midrule
CRRA & $B=0$ & 递减、双曲线 & 风险金比例不随财富变 \\
CARA & $A=0$ & 常数 & 风险金绝对额不随财富变 \\
二次型 & $A<0$ & 递增 & 有餍足点，越富越避险 \\
一般 HARA & $A>0,B>0$ & 递减 & 含“基本需求”式底线 $B$ \\
\bottomrule
\end{tabular}
\end{center}
统一形式 $u(w)=\dfrac{1-\gamma}{\gamma}\Bigl(\dfrac{aw}{1-\gamma}+b\Bigr)^{\gamma}$。
HARA 的可贵之处：组合分离（两基金分离定理）在该族成立。

\subsection{状态依赖效用与失望厌恶}

\notehead{状态依赖}
$u=u(c,s)$：效用直接依赖状态 $s$（健康、生存状态）。
例：健康状态下钱的边际效用高于患病状态。
后果：最优保险不再是全额保险——即使保费公平，
也可能只部分投保（状态间边际替代率不再等于 1）。
适用：健康经济学、寿险、长护险。

\notehead{失望厌恶}
\[
U=\mathbb{E}[u(c)]-\theta\,\mathbb{E}\bigl[\max\{\mathcal{CE}-c,\,0\}\bigr],
\]
以确定性等价 $\mathcal{CE}$ 为内生参照点：结果低于 $\mathcal{CE}$ 的部分
被额外加权 $\theta$。\textbf{一阶风险厌恶}：
对小赌局的风险金与赌局规模同阶（期望效用下为二阶），
可同时解释“对小赌过度规避”与“股权溢价”，而不必把 $\gamma$ 设得离谱。

\notehead{幂指数型}
$u(w)=-\exp(-\alpha w^{\theta})$（$\alpha,\theta>0$）：
两参数使 ARA 形态更灵活（$\theta<1$ 递减 ARA + 递减 RRA 的组合），
用于拟合实验数据中风险态度随财富的变化。

% ============================================================
\section{模糊、非期望效用与自控（补全）}
% ============================================================

\subsection{最大最小期望效用（多先验/模糊厌恶）}\label{subsec:meu}

\notehead{动机：埃尔斯伯格悖论}
一缸 90 球：红 30，黑与黄各 0--60（比例未知）。
赌局 A：抽红得 100；赌局 B：抽黑得 100。
多数人选 A（已知概率）而非 B（未知概率）；
但对“红或黄得 100”与“黑或黄得 100”，同一批人又选后者。
两对选择合在一起\textbf{没有任何单一概率分布}能解释：
若 $p(\text{红})>p(\text{黑})$（第一对），则 $p(\text{黑或黄})>p(\text{红或黄})$（第二对）——矛盾。
期望效用的“信念必为单一概率”假设被证伪。

\notehead{形式}
\[
U(f)=\min_{p\in\mathcal{C}}\ \mathbb{E}_p\,u(f),
\]
$\mathcal{C}$ 为先验集合（对未知成分的概率“只确定一个范围”）。
模糊厌恶 = 按\textbf{最坏}先评估。上例中
$\mathcal{C}=\{p:\ p(\text{红})=1/3,\ p(\text{黑})\in[0,2/3]\}$，
A 的期望确定（$u(100)/3$），B 的最坏期望为 $0$——选 A；
“黑或黄”得 100：$p(\text{黑})+p(\text{黄})=2/3$ 对一切先验恒定，期望 $=\frac23u(100)$ 无模糊；
“红或黄”得 100：$p(\text{红或黄})=\frac13+p(\text{黄})\in[\frac13,1]$，最坏期望 $=\frac13u(100)$——
故选“黑或黄”。悖论被容纳。

\notehead{性质与变体}
\begin{itemize}
  \item 无差异曲线在确定性线（45° 线）处出现\textbf{折点}（一阶风险厌恶的来源之一）；
  \item $\alpha$-最大最小：$U(f)=\alpha\min_{\mathcal{C}}\mathbb{E}u+(1-\alpha)\max_{\mathcal{C}}\mathbb{E}u$，
        $\alpha$ 度量悲观程度，可容纳模糊中性/爱好；
  \item 平滑模糊厌恶：$U=\phi^{-1}\Bigl(\mathbb{E}_{\mu}\,\phi\bigl(\mathbb{E}_p u\bigr)\Bigr)$，
        二阶信念 $\mu$（对先验的先验）经凹函数 $\phi$ 聚合——
        把“对风险的态度”与“对模糊的态度”分置两层，便于资产定价应用。
\end{itemize}

\begin{center}
\begin{tikzpicture}[scale=1.0]
  \draw[axisline] (0,0) -- (5.6,0) node[below left, axislabel] {状态 1 消费 $c_1$};
  \draw[axisline] (0,0) -- (0,4.6) node[below left, axislabel] {状态 2 消费 $c_2$};
  \draw[dashed, MidBlue] (0,0) -- (3.4,3.4);
  \node[below right, align=center] at (1.15,3.75) {确定性线};
  % kinked ICs at 45 line
  \draw[thick, MidBlue] (0.6,3.9) -- (2.1,2.1) -- (4.4,1.4);
  \draw[thick, MidBlue] (0.6,2.7) -- (1.4,1.4) -- (4.4,0.35);
  \node[align=left] at (4.6,3.3) {多先验下的\\ 无差异曲线：\\ 折点在确定性线上\\ （对两状态的概率\\ 各按最坏情形评估）};
\end{tikzpicture}
\end{center}
\figexplan{横纵轴为两状态下的消费。虚线为确定性线（两状态消费相等）；无差异曲线在确定性线上出现折点，因为多先验偏好对两个状态分别按各自的最坏先验评估，折点两侧的有效概率不同。折点恰落在确定性线上，是这类偏好区别于期望效用的几何特征。}

\notehead{为什么选它}
宏观金融（政策规则的模糊、模型不确定性下的稳健决策）、
气候变化（概率本身未知的深不确定性）、健康冲击。
选择标准：若“连概率都不知道”是问题的核心（而非修辞），
期望效用不是近似而是\textbf{误设}——必须用多先验类形式。

\subsection{秩依赖效用：概率加权}

\notehead{形式与推导}
把结果从低到高排序 $x_1\le x_2\le\cdots\le x_n$，权重
\[
\pi_i=w\Bigl(\sum_{j\ge i}p_j\Bigr)-w\Bigl(\sum_{j>i}p_j\Bigr),
\qquad U=\sum_i\pi_i u(x_i),
\]
$w:[0,1]\to[0,1]$ 递增、$w(0)=0,w(1)=1$ 为\textbf{概率加权函数}。
$\pi_i$ 是“结果 $x_i$ 在累积分布尾部增量上的权重”：
$w$ 的曲率决定权重如何向极端结果倾斜。
典型反 S 形（小概率放大、大概率压缩）：买彩票（放大尾部小概率收益）
与买保险（放大尾部小概率损失）在同一人身上并存——
期望效用需要 $u$ 既凸又凹才能解释，秩依赖用一个 $w$ 即可。

\begin{center}
\begin{tikzpicture}[scale=0.95]
  \draw[axisline] (0,0) -- (5.0,0) node[below left, axislabel] {累积概率 $p$};
  \draw[axisline] (0,0) -- (0,4.6) node[below left, axislabel] {$w(p)$};
  \draw[dashed, MidBlue] (0,0) -- (4.2,4.2);
  \draw[thick, MidBlue, domain=0:4.2, samples=90] plot ({\x},{4.2*((\x/4.2)^0.61)/(((\x/4.2)^0.61+(1-\x/4.2)^0.61)^(1/0.61))});
  \node[align=left] at (3.0,1.2) {反 S 形：\\ 小概率放大\\ 大概率压缩};
\end{tikzpicture}
\end{center}
\figexplan{横轴为累积概率 $p$，纵轴为加权后的决策权重 $w(p)$。虚线是 $w(p)=p$ 的参照线（期望效用）；实线呈反 S 形：小概率被放大（左下高于虚线）、大概率被压缩（右上低于虚线）。同一条曲线因此可同时解释买彩票与买保险。}

\notehead{为什么选它}
保留序的单调性（一阶随机占优不被违反——前景理论的原始加权函数做不到），
又能解释共同比率效应、阿莱悖论、彩票与保险并存。
是行为风险模型中\textbf{理论与数据兼容度最高}的一支。

\subsection{后悔理论}

\notehead{形式}
两选项 $A,B$ 的结果对 $(x^A_s,x^B_s)$ 在状态 $s$ 实现。
效用含\textbf{比较项}：
\[
U(A)=\sum_s p_s\Bigl[u(x^A_s)+Q\bigl(u(x^A_s)-u(x^B_s)\bigr)\Bigr],
\]
$Q$ 递增且 $Q(0)=0$（$Q$ 凹于负域时产生后悔厌恶）。
选择 $A$ 后若 $B$ 的结果更好，承受“本可以”的负效用。

\notehead{能反映什么}
偏好反转（§4.3 第三章）、事后归因改变事前选择
（因担心踏空而跟风投资、因担心后悔而过度投保）。
与前景理论的区别：参照点不是状态变量，而是\textbf{未选项的反事实}。

\subsection{诱惑与自控：菜单上的偏好}

\notehead{形式（Gul--Pesendorfer）}
偏好定义在\textbf{菜单}（选择集）上，效用
\[
U(A)=\max_{x\in A}\,\bigl[u(x)+v(x)\bigr]-\max_{y\in A}\,v(y),
\]
$u$：长期（规范性）偏好；$v$：诱惑（短期）偏好。
第二项是\textbf{自控成本}：菜单里最有诱惑力的选项即使没选，
也消耗自控资源。

\notehead{数值例}
沙拉：$u=3,v=0$；蛋糕：$u=0,v=2$。
\begin{itemize}
  \item 菜单 $\{\text{沙拉}\}$：$U=3-0=3$；
  \item 菜单 $\{\text{沙拉,蛋糕}\}$：$\max(u+v)=3$（选沙拉），$U=3-2=1$。
\end{itemize}
$U(\{\text{沙拉}\})>U(\{\text{沙拉,蛋糕}\})$：\textbf{菜单越大越差}——
理性框架中不可能（选择集扩大弱占优），
它刻画的是\textbf{承诺需求}：主动删除选项（注销账户、不囤零食）。
若蛋糕 $v=4$：选蛋糕（$u+v=4>3$），$U=4-4=0$——抵抗失败，效用归零。

\begin{center}
\begin{tikzpicture}[scale=1.0, font=\small]
  \node[draw=AxisBlue, rounded corners=2pt, align=center, fill=MidBlue!10] (m1) at (0,0) {$\{\text{沙拉}\}$\\ $U=3$};
  \node[draw=AxisBlue, rounded corners=2pt, align=center, fill=MidBlue!10] (m2) at (4.2,0) {$\{\text{沙拉,蛋糕}\}$\\ $U=1$（抵抗成功）};
  \node[draw=AxisBlue, rounded corners=2pt, align=center, fill=MidBlue!10] (m3) at (9.2,0) {$\{\text{沙拉,蛋糕}^*\}$\\ $U=0$（抵抗失败）};
  \draw[-{Stealth}, thick, color=AxisBlue] (m1) -- (m2) node[midway, above] {加入诱惑};
  \draw[-{Stealth}, thick, color=AxisBlue] (m2) -- (m3) node[midway, above] {诱惑增强 $v=2\to4$};
  \node[align=center] at (4.6,-1.5) {偏好承诺：宁可面对更小的菜单——选择集扩大不再弱占优};
\end{tikzpicture}
\end{center}
\figexplan{三个方框是三个可选菜单及其承诺效用：只有沙拉时效用为 $3$；加入蛋糕后虽可抵抗，但要付出自控成本，效用降为 $1$；诱惑进一步增强（$v=2\to4$）时抵抗失败，效用归零。箭头表示菜单扩大与诱惑增强的方向——选择集扩大不再弱占优，这正是承诺需求的含义。}

\notehead{为什么选它}
与 $\beta$-$\delta$（§4.3.4）是自控问题的两种语法：
$\beta$-$\delta$ 用时间结构（计划反转），
本形式用偏好结构（承诺偏好、自控成本）。
后者可直接做福利分析：自控成本有明确的效用单位，
节食计划、储蓄承诺产品的定价都可模型化。

\subsection{参照依赖的获得—损失效用}

\notehead{形式（Kőszegi--Rabin）}
\[
u(c\mid r)=m(c)+\mu\bigl(m(c)-m(r)\bigr),\qquad
\mu(z)=
\begin{cases}
\eta z, & z\ge0,\\
\eta\lambda z, & z<0,
\end{cases}
\quad \lambda>1,
\]
$m(c)$：消费效用；$\mu$：相对参照点 $r$ 的获得—损失效用，
$\eta$ 为获得—损失效用的权重。
关键创新：\textbf{参照点 $r$ 内生于期望}——
$ r$ 是个体事前对结果的理性预期分布（个人均衡：
预期一致，即实际选择所隐含的分布正是参照点）。

\notehead{能反映什么（统一账户）}
禀赋效应（拥有后参照点含该物，卖出是损失）、
消费对收入冲击的过度敏感（预期调整滞后时，意外收入是获得）、
劳动供给的目标收入行为、拍卖中的虚高报价。
它把第三章 §4.8 罗列的一批现象压缩为\textbf{一个两参数结构}（$\eta,\lambda$），
且保持期望效用的大部分可操作性——这是它成为应用基准的原因。

\subsection{心理账户与选择括定}

\notehead{形式}
不是单一效用函数，而是\textbf{评价的结构}：
总决策被拆分为子账户（日常开销/旅行基金/养老储蓄），
各账户有各自的效用与预算约束，账户间资金不完全可替代；
宽括定 = 单一账户联合最大化；窄括定 = 逐账户独立最大化。

\notehead{能反映什么}
边际消费倾向随资金来源变化（意外之财花得快、工资存得多）、
“ sunk cost ”支付后继续消费、养老金自动加入的高留存率
（默认选项被归入“不可动账户”）。
§4.7.4 的 AER 2024 文献给出从选择数据识别括定方式的公理化检验：
先实证判断当事人如何记账，再决定效用该写成一个还是一叠。

% ============================================================
\section{行为效用、社会偏好与离散选择}
% ============================================================

\subsection{参照依赖价值函数（前景理论型）}

\notehead{形式与性质}
\[
v(z)=
\begin{cases}
z^{\alpha}, & z\ge0,\\[2pt]
-\lambda(-z)^{\beta}, & z<0,
\end{cases}
\qquad \alpha,\beta\in(0,1),\ \lambda>1.
\]
三个结构性特征，各对应一类实验事实：
\begin{enumerate}
  \item \textbf{参照依赖}：自变量是\textbf{相对参照点的损益} $z$，不是财富水平；
  \item \textbf{损失厌恶}：$\lambda>1$，损失的效用权重是等额收益的 $\lambda$ 倍
        （估计值常落在 $2$ 附近）；
  \item \textbf{敏感度递减}：$\alpha,\beta<1$，损益越大边际感受越小——
        价值函数在收益域凹、损失域凸（S 形）。
\end{enumerate}
直接推论：\textbf{反射效应}——收益域风险规避、损失域风险偏好
（“$+100$ 确定”对“$+200$ 半概率”，多数人选前者；
“$-100$ 确定”对“$-200$ 半概率”，多数人选后者）。
期望效用无法在一个 $u$ 下同时产生这两个选择。

\begin{center}
\begin{tikzpicture}[scale=0.95]
  \draw[axisline] (-3.4,0) -- (3.6,0) node[below left, axislabel] {$z$（相对参照点）};
  \draw[axisline] (0,-3.0) -- (0,2.6) node[below left, axislabel] {$v(z)$};
  \draw[thick, MidBlue, domain=0:3.2, smooth, variable=\x] plot ({\x},{sqrt(\x)*1.15});
  \draw[thick, MidBlue, domain=0:3.2, smooth, variable=\x] plot ({-\x},{-(\x)^0.6*1.55});
  \node[nodeblue] at (0,0) {};
  \draw[dashed, MidBlue] (0,0) -- (2.9,2.35) node[right, align=left] {损失域切线更陡：\\ 斜率比 $\approx\lambda$};
  \node[align=left] at (2.3,-1.5) {原点折点（损失厌恶）\\ 收益凹、损失凸\\ （敏感度递减）};
\end{tikzpicture}
\end{center}
\figexplan{横轴为相对参照点的损益 $z$，纵轴为价值 $v(z)$。曲线在原点出现折点：收益域为凹（风险厌恶）、损失域为凸（风险寻求），且损失域切线更陡（虚线），斜率之比约为损失厌恶系数 $\lambda$。原点处的折点即参照依赖，两侧曲率递减对应敏感度递减。}

\notehead{适用情形}
资产配置（处置效应：卖盈持亏）、劳动供给（日收入目标）、
定价与营销（免费试用后的退订损失）、保单设计。
与概率权重函数配合即为完整前景理论；仅价值函数已能解释大多数参照点现象。

\subsection{不平等厌恶（收益分配型社会偏好）}\label{subsec:fs}

\notehead{形式一：差异厌恶}
\[
u_i=x_i-\alpha_i\frac{1}{n-1}\sum_{j\ne i}\max\{x_j-x_i,\,0\}
-\beta_i\frac{1}{n-1}\sum_{j\ne i}\max\{x_i-x_j,\,0\},
\qquad 0\le\beta_i\le\alpha_i.
\]
$\alpha_i$：对\textbf{不如人}的厌恶（劣势不平等）；
$\beta_i$：对\textbf{优于人}的不安（优势不平等），且设定 $\beta_i\le\alpha_i$
——不如人的痛苦强于优于人的不安，这是实验数据的稳健结论。
两人情形简化为
\[
u_i=
\begin{cases}
x_i-\beta_i(x_i-x_j), & x_i\ge x_j,\\[2pt]
x_i-\alpha_i(x_j-x_i), & x_i<x_j.
\end{cases}
\]
无差异曲线在 $x_i=x_j$ 处成折线，折点两侧斜率不同（$\alpha$ 侧更陡）。

\begin{center}
\begin{tikzpicture}[scale=1.0]
  \draw[axisline] (0,0) -- (5.6,0) node[below left, axislabel] {$x_i$};
  \draw[axisline] (0,0) -- (0,4.4) node[below left, axislabel] {$x_j$};
  \draw[dashed, MidBlue] (0,0) -- (4.2,4.2) node[above right] {$x_i=x_j$};
  % kinked ICs: lines with kink on diagonal
  \draw[thick, MidBlue] (0.5,3.8) -- (2.2,2.2) -- (4.9,1.5);
  \draw[thick, MidBlue] (0.5,2.6) -- (1.4,1.4) -- (4.9,0.35);
  \node[align=left] at (4.85,3.05) {折线无差异曲线：\\ 对角线左侧（不如人）\\ 斜率更陡（$\alpha$ 侧）};
\end{tikzpicture}
\end{center}
\figexplan{虚线为 $x_i=x_j$ 的平等线。无差异曲线在该线上折曲：对角线左侧（自己不如他人）斜率更陡，对应不利不平等厌恶参数 $\alpha$；右侧斜率较缓，对应有利不平等厌恶 $\beta$。折点两侧斜率不等，正是分配偏好无法用光滑效用表示的表现。}

\notehead{形式二：社会最优与最小者关心}
\[
u_i=(1-\rho-\sigma)x_i+\rho\min(x_i,x_j)+\sigma\max(x_i,x_j).
\]
$\rho>0$ 时在意\textbf{最差者}（慈善/效率兼顾取决于 $\sigma$ 符号与大小）；
$\sigma<0$ 时甚至\textbf{嫉妒}（宁损己以损人）。
与差异厌恶形式的区别：本形式关心极端值而非逐对差距，
对分配中的“最底层”更敏感。

\notehead{适用情形}
最后通牒/独裁者/公共品博弈的结构估计、
组织内薪酬压缩的福利分析、再分配偏好的调查实验。
估计要点：$\alpha$ 由拒绝行为识别，$\beta$ 由分配行为识别。

\subsection{离散选择：随机效用模型}\label{subsec:rum}

\notehead{形式与 Logit 推导}
个体 $i$ 对方案 $j$ 的效用
\[
u_{ij}=v_j(\boldsymbol{x}_j;\theta)+\varepsilon_{ij},
\]
$v_j$ 为系统部分（可观测属性的函数），$\varepsilon_{ij}$ 为研究者不可观测的口味扰动。
设 $\{\varepsilon_{ij}\}$ 独立同分布于 I 型极值分布 $F(\varepsilon)=\exp(-e^{-\varepsilon})$，则
\[
\begin{aligned}
P_i(j)&=\Pr\bigl(v_j+\varepsilon_{ij}\ge v_k+\varepsilon_{ik},\ \forall k\bigr)\\
&=\int\prod_{k\ne j}F\bigl(v_j-v_k+\varepsilon\bigr)\,f(\varepsilon)\,\mathrm{d}\varepsilon
=\boxed{\ \frac{e^{v_j}}{\sum_{k}e^{v_k}}\ }.
\end{aligned}
\]
关键一步：I 型极值分布的稳定性——独立极值变量之差服从 Logistic 分布，
多重积分闭式可解。选择概率只依赖\textbf{效用差}，与绝对水平无关。

\notehead{性质与扩展}
\begin{itemize}
  \item \textbf{IIA 性质}：两方案的选择概率比 $P(j)/P(k)=e^{v_j-v_k}$ 与第三方无关——
        既是计算优点也是经验缺陷（红蓝巴士悖论）；
  \item \textbf{嵌套 Logit}：方案分组（巢内相关），放松 IIA——
        与 §\ref{subsec:flexible} 分支树型效用互为表里；
  \item \textbf{随机系数 Logit}：$\theta$ 在个体间随机分布，
        允许替代模式任意化，是需求估计的工业标准（§\ref{subsec:jpe-markups}）。
\end{itemize}

\notehead{适用情形}
商品选择（品牌、车型）、区位选择、职业与学校选择、
任何“多选一”且需要\textbf{从选择数据反推偏好}的经验场景。
福利计算有闭式公式：期望最大效用（消费者剩余）
$\mathbb{E}\max_j u_{ij}=\ln\sum_j e^{v_j}+$ 常数（对数和公式）。

\subsection{类型学小结：六个选择维度}

\begin{center}
\small
\begin{tabular}{llll}
\toprule
维度 & 问题 & 形式分支 & 代表 \\
\midrule
商品结构 & 替代还是互补？ & 线性 / 乘积型 / CES / 固定比例 & §4.1 \\
收入效应 & 份额恒定还是变动？ & 位似 / 基本需求型 / 非位似 & §4.1.5--§4.1.7, §4.2 \\
财富效应 & 是否可忽略？ & 拟线性 vs 一般 & §4.1.6 \\
时间 & 一致还是不一致？ & 指数贴现 / $\beta$-$\delta$ & §4.3.1, §4.3.4 \\
风险 & 时间可加还是解绑？ & CRRA / 递归型 / 失望厌恶 & §4.3.3, §4.3.5--§4.3.6 \\
他人 & 自利的还是社会的？ & 自利 / 差异厌恶 / 攀比 & §4.3.2, §4.5.2 \\
\bottomrule
\end{tabular}
\end{center}

% ============================================================
\section{社会福利函数：从个人效用加总成社会评价（补全）}
% ============================================================

政策评估需要把 $N$ 个人的效用（或收入）映射为一个数——
这个映射本身就是一类“效用函数”，其形式选择\textbf{直接编码分配价值观}。

\subsection{功利主义（效用加总型）}

\notehead{形式}
$W=\sum_i u_i$。

\notehead{性质}
每个人的效用权重相同；对分配\textbf{完全中性}：
一元从穷人转给富人，$W$ 不变（若边际效用不被考虑）；
若 $u$ 凹，则同样总量下平均分配最优——\textbf{分配关怀全部来自个体效用曲率，
而非 $W$ 的形式}。
优点：与补偿原则（总剩余最大化）一致，是成本收益分析的默认件；
缺点：无法接受“一元对穷人的价值更高”这类独立的分配判断。

\begin{figure}[H]
\centering
\begin{tikzpicture}[font=\scriptsize,scale=0.88,transform shape]
  \def\frontier{(0,0.6) .. controls (2.4,3.3) and (1.5,3.5) .. (3.3,3.6)}
  % (a) 功利主义
  \draw[axisline] (0,0) -- (3.45,0) node[below right] {$u_1$};
  \draw[axisline] (0,0) -- (0,3.9) node[above left] {$u_2$};
  \draw[thick] \frontier;
  \draw[dashed,thin] (0,3.6) -- (3.3,3.6);
  \draw[->,thick] (0,3.6) -- (3.3,0) node[midway,above left,font=\scriptsize] {$u_1+u_2=\text{常数}$};
  \fill (1.95,1.65) circle (1.5pt);
  \node[font=\scriptsize, align=center] at (1.95,1.20) {(a) 功利主义：\\最大化总和};
  % (b) 最大最小
  \begin{scope}[xshift=5.3cm]
    \draw[axisline] (0,0) -- (3.45,0) node[below right] {$u_1$};
    \draw[axisline] (0,0) -- (0,3.9) node[above left] {$u_2$};
    \draw[thick] (0,0.6) .. controls (2.4,3.3) and (1.5,3.5) .. (3.3,3.6);
    \draw[dashed,thin] (0.0,2.78) -- (2.42,2.78);
    \draw[->,thick] (0,2.78) -- (2.58,0) node[midway,above left,font=\scriptsize] {$u_2=\text{常数}$};
    \fill (2.42,2.78) circle (1.5pt);
    \node[align=center,font=\scriptsize] at (2.05,1.35) {(b) 最大最小：\\抬高最差者的水平};
  \end{scope}
  % (c) CES
  \begin{scope}[xshift=10.6cm]
    \draw[axisline] (0,0) -- (3.45,0) node[below right] {$u_1$};
    \draw[axisline] (0,0) -- (0,3.9) node[above left] {$u_2$};
    \draw[thick] (0,0.6) .. controls (2.4,3.3) and (1.5,3.5) .. (3.3,3.6);
    \draw[->,thick] (0,2.20) to[bend left=22] (2.95,0.66) node[midway,above,font=\scriptsize] {等 $\varepsilon$ 无差异曲线};
    \fill (2.16,1.20) circle (1.5pt);
    \node[align=center,font=\scriptsize] at (2.05,2.42) {(c) CES：$\varepsilon$ 居中，\\切于两极限之间};
  \end{scope}
\end{tikzpicture}
\caption{三种社会福利函数在同一可能性前沿上的最优分配}
\end{figure}
\figexplan{三幅图共用同一条效用可能性前沿（实线），区别只在等社会福利线。功利主义用斜率为 $-1$ 的直线（$u_1+u_2$ 为常数），最优点在两效用较均衡处；最大最小用水平线（$u_2$ 为常数），最优点是把最差者的效用推到最高；CES 型用一族曲率居中的无差异曲线，切点介于前两者之间。图形要点：福利函数的选择等价于在可能性前沿上“选点”，参数 $\varepsilon$ 正是这一选择的调节旋钮。}

\subsection{最大最小（最差者优先型）}

\notehead{形式}
$W=\min_i u_i$。

\notehead{性质}
社会福利等于\textbf{最差者}的效用；任何不改善最差者的再分配都无价值。
无差异曲线为 L 形（与个人效用的完全互补同构，§4.1.3）：
只有最差者改善，$W$ 才改善。
极限性质：它对不平等“无限厌恶”——宁可总量大幅缩水也要托底。

\subsection{CES 型社会福利函数（含不平等厌恶参数）}

\notehead{形式}
\[
W(x)=\Bigl[\frac1N\sum_i x_i^{\,1-\varepsilon}\Bigr]^{\frac{1}{1-\varepsilon}},
\qquad \varepsilon\ge0,\ \varepsilon\ne1,
\]
$\varepsilon$ 为\textbf{不平等厌恶}参数；$\varepsilon=1$ 时取几何平均
$\prod_i x_i^{1/N}$。

\notehead{两个极限（推导）}
$\varepsilon\to0$：$W\to\frac1N\sum_i x_i$——功利主义（算术平均）；
$\varepsilon\to\infty$：幂均值的极限为最小值，$W\to\min_i x_i$——最大最小。
故 CES 型是\textbf{功利主义与最大最小之间的连续桥梁}，
$\varepsilon$ 就是这座桥上的刻度。

\notehead{均等分布等价与阿特金森指数}
定义\textbf{均等分布等价}（EDE）$\xi=W(x)$：
若每人得 $\xi$，社会福利与现实分配相同。不平等指数
\[
A=1-\frac{\xi}{\bar x},\qquad \bar x=\frac1N\sum_ix_i.
\]
\textbf{数值例}：两人收入 $(10,2)$，$\bar x=6$。
\begin{center}
\small
\begin{tabular}{lccc}
\toprule
$\varepsilon$ & $\xi$（计算） & $A=1-\xi/6$ & 含义 \\
\midrule
$0$（功利主义） & $6.0000$ & $0.0000$ & 只看总量 \\
$1$（几何平均） & $\sqrt{20}=4.4721$ & $0.2546$ & 温和厌恶 \\
$2$ & $(0.5(10^{-1}+2^{-1}))^{-1}=3.3333$ & $0.4444$ & 强厌恶 \\
$\to\infty$（最大最小） & $\to2.0000$ & $\to0.6667$ & 只看最差者 \\
\bottomrule
\end{tabular}
\end{center}
（复算 $\varepsilon=2$：$W=[0.5(10^{-1}+2^{-1})]^{-1}=[0.5\times0.6]^{-1}=1/0.3=3.3333$。）

\begin{center}
\begin{tikzpicture}[scale=0.95]
  \draw[axisline] (0,0) -- (5.9,0) node[below left, axislabel] {$u_1$};
  \draw[axisline] (0,0) -- (0,4.4) node[below left, axislabel] {$u_2$};
  % epsilon=0: straight line
  \draw[thick, MidBlue] (0,3.4) -- (4.4,1.0);
  \node[right] at (4.4,1.05) {$\varepsilon=0$};
  % epsilon=1: hyperbola
  \draw[thick, MidBlue, domain=0.8:5.0, smooth, variable=\x] plot ({\x},{2.2/\x});
  \node[above right] at (3.55,1.05) {$\varepsilon=1$};
  % epsilon large: near L
  \draw[thick, MidBlue] (0.95,3.8) .. controls (1.0,1.45) and (1.4,1.0) .. (4.6,0.75);
  \node[right] at (4.6,0.62) {$\varepsilon$ 大};
  \node[align=left] at (2.9,3.8) {社会无差异曲线：\\ 曲率 $=$ 不平等厌恶程度\\ L 形极限 $=$ 最大最小};
\end{tikzpicture}
\end{center}
\figexplan{横纵轴为两人的效用。$\varepsilon=0$ 时为直线（功利主义，只关心总和）；$\varepsilon=1$ 时为双曲线（乘法型）；$\varepsilon$ 越大越接近 L 形，极限即最大最小（罗尔斯）准则。曲率因此可直接读作不平等厌恶程度：政策评估只需随 $\varepsilon$ 重算一遍。}

\notehead{为什么选它}
(i) 一参数覆盖全部主流分配立场，政策报告的\textbf{分配敏感性}
就是随 $\varepsilon$ 重算一遍；
(ii) EDE 与不平等指数有闭式，沟通成本低；
(iii) 与个体 CES 同构，宏观模型中“社会规划者偏好”可直接借用消费品加总技术。

\subsection{Kolm--Pollak 型：绝对不等价}

\notehead{形式}
\[
W(x)=-\frac{1}{\kappa}\ln\Bigl(\frac1N\sum_i e^{-\kappa x_i}\Bigr),
\qquad \kappa>0,
\]
\textbf{性质}：给所有人同时增加等量收入，$W$ 增加相同的量
（平移不变性）——不平等按其\textbf{绝对差距}度量；
CES 型则按\textbf{比例差距}（$x$ 与 $2x$ 分布被认为同样不平等）。
选择标准：认为“贫富差距一万”与“贫富差距十万，但同比例”一样不平等
→ Kolm--Pollak；认为不同 → CES 型。
$W$ 同时是\textbf{确定性等价}结构：与 CARA 同族，
税率/转移的福利分析中常与 CRRA 个体效用搭配使用。

\subsection{几何平均与议价型福利函数}

\notehead{形式}
$W=\prod_i u_i^{1/N}$（CES 型 $\varepsilon=1$ 特例），
等价于 $\ln W=\frac1N\sum_i\ln u_i$。
两大用途：
(i) 两人、可变权重时即为议价解的目标函数（§4.2.6），
\textbf{社会选择与合作解在此汇合}；
(ii) 代际福利：$W=\prod_t c_t^{\beta_t}$ 对不完全替代的不同时期消费
赋予零下限（任何一期消费为零，$W=0$）——
气候经济学中“不允许未来被完全牺牲”的形式化方式之一。

\subsection{社会福利函数选择对照表}

\begin{center}
\small
\begin{tabular}{llll}
\toprule
形式 & 不平等态度 & 关键性质 & 典型用途 \\
\midrule
功利主义 $\sum u_i$ & 中性 & 与补偿原则一致 & 成本收益分析 \\
最大最小 $\min u_i$ & 无限厌恶 & 只认最差者 & 托底政策、正义论 \\
CES 型（$\varepsilon$） & 一参数可调 & 比例差距、EDE 闭式 & 分配敏感性报告 \\
Kolm--Pollak（$\kappa$） & 一参数可调 & 绝对差距、平移不变 & 税与转移分析 \\
几何平均 & 温和厌恶 & 与议价解同构 & 代际福利 \\
\bottomrule
\end{tabular}
\end{center}

% ============================================================
\section{顶刊前沿：近年 AER/JPE 怎样用效用函数（深化版）}
% ============================================================

本节基于 2021--2026 年 AER、JPE 及相关顶刊的公开文献（均经联网核实），
按 \textbf{文献 → 模型（公式 + 假设）→ 新质 → 可借鉴处} 拆解。
重点不在“他们用了哪个名字的形式”，而在\textbf{效用函数的自变量、
参数化与识别策略}如何被重新设计。

\subsection{随机系数 Logit 的偏好面板化：D\"opper--MacKay--Miller--Stiebale（JPE 2025）}\label{subsec:jpe-markups}

\notehead{模型}
个体 $i$ 在市场 $t$（品类 $\times$ 年）购买产品 $j$ 的间接效用：
\[
u_{ijt}=\underbrace{x_{jt}\beta-\alpha p_{jt}+\xi_{jt}}_{\delta_{jt}\ \text{(均值效用)}}
+\underbrace{\sum_k x_{jkt}\,\sigma_k\,\nu_{ik}}_{\mu_{ijt}\ \text{(个体异质)}}
+\varepsilon_{ijt},
\]
其中 $\nu_{ik}\sim N(0,1)$ 为个体口味冲击，$\sigma_k$ 度量偏好的离散程度，
$\varepsilon_{ijt}$ 为 I 型极值扰动（§\ref{subsec:rum}）。

\notehead{关键假设}
(i) $\mathbb{E}[\xi_{jt}\mid Z_{jt}]=0$：工具变量（成本、竞品特征）与不可观测质量正交；
(ii) 企业按 Nash--Bertrand 定价：价格加成 $=$ 需求弹性倒数（一阶条件），
从而\textbf{边际成本可由需求参数反推}；
(iii) 偏好在品类内随年份变化——每个品类 $\times$ 年\textbf{单独}估计
$(\alpha_t,\sigma_t,\delta_{jt})$。

\notehead{新质}
传统用法把 $(\alpha,\sigma)$ 当跨期结构常数；本文得到\textbf{偏好面板}：
2006--2019 年 $\alpha_t$（价格敏感度）系统性下降，
对应消费者剩余上移与企业加成率平均上升约 30\%。
偏好参数从“待估常数”变成“待解释的时变对象”——
这是需求估计范式的转向。

\notehead{可借鉴}
任何长期政策评估（养老金、医疗市场）若样本跨十年以上，
“偏好稳定”假设应先被检验而非默认；
检验成本不高：分时段重估并比较参数即可。

\subsection{福利分析对函数形式的稳健性：Kang--Vasserman（AER 2025）}

\notehead{模型}
不指定需求函数形式。设真实需求 $q(p)$ 未知，
分析者假设某形式 $\hat q(p)$ 并据此算福利 $\hat W$。
定义\textbf{稳健半径}：在梯度与曲率允许偏差
$\|\nabla q-\nabla\hat q\|\le\Delta_1$、$\|\nabla^2 q-\nabla^2\hat q\|\le\Delta_2$
内，使福利结论反转所需的\textbf{最小} $(\Delta_1,\Delta_2)$。
用信息设计工具把最小化问题化为可计算规划。

\notehead{新质}
把“换两组参数重算”的朴素稳健性升级为\textbf{形式无关的度量}：
回答的不是“参数错了怎么办”，而是“形式错了多远结论才翻”。

\notehead{可借鉴}
政策论文的福利部分可附一行稳健半径：
“需求曲率需偏离假设形式 X\% 以上，本结论才反转”——
比传统敏感性表格信息密度高一个量级（§4.8 第 9 步已纳入）。

\subsection{非位似 CES 与结构转型：Comin--Lashkari--Mestieri（ECMA 2021 及 2025 后续）}

\notehead{模型}
偏好由隐式加法可分形式定义（§4.2.3）：
$\sum_i\Omega_i(c_iU^{\varepsilon_i})^{(\sigma-1)/\sigma}=1$。
关键可估对象：相对恩格尔曲线斜率 $\varepsilon_i-\varepsilon_j$，
用美国家庭、印度家庭及 39 国历史面板分别估计。

\notehead{关键假设}
(i) 部门收入弹性 $\varepsilon_i$ 为常数（恩格尔斜率不衰减）；
(ii) 价格效应（$\sigma$）与收入效应（$\varepsilon_i$）分离进入。

\notehead{新质与结论}
收入效应解释了\textbf{国家内部}结构转型的大部分；
位似偏好（$\varepsilon_i\equiv1$）在数据上被明确拒绝。
2025 年后续工作把同一偏好嵌入创新模型：
收入弹性高的部门（服务）吸引更多研发——
需求结构\textbf{引导}技术进步方向。

\notehead{可借鉴}
研究中国消费升级/产业升级时，位似 CES 的“份额恒定”
假设直接排除核心机制；非位似 CES 是当前可估计的标准件。

\subsection{选择括定的显示偏好检验：Ellis--Freeman（AER 2024）}

\notehead{模型}
观察个体在一系列\textbf{同时呈现的}预算集上的选择。
若个体宽括定：所有预算的联合约束下最大化单一效用；
若窄括定：每个预算独立最大化各自的（可能不同的）效用。
文章给出两组选择数据上的公理（弱/强括定的一致性条件），
无需指定效用函数参数形式即可判定括定类型。

\notehead{新质}
把“心理账户”从实验设定变量变为\textbf{可从常规选择数据识别}的对象；
方法完全非参数（与本课程第二章显示偏好路线同源）。

\notehead{可借鉴}
养老金、医疗共付、教育储蓄等“天然多账户”场景，
建模前应先问括定方式；宽/窄括定下同一政策的福利评估可以反向。

\subsection{随机效用的可观测性边界：Kono--Saito--Sandroni（AER 2026）}

\notehead{模型}
标准 RUM（§\ref{subsec:rum}）假设所有备选的选择频率可观测；
现实中“外部选项”（未购买、未申报）的频率常缺失。
文章推导：当部分方案频率不可观测时，
刻画 RUM 可理性化的\textbf{有限、非冗余不等式组}。
核心发现：把不可观测方案笼统聚合为一个“外部选项”
会丢失 RUM 的可检验内容——\textbf{即使数据与聚合后的 RUM 一致，
也可能不存在任何与原始数据一致的随机效用模型}。

\notehead{可借鉴}
用 Logit 族估计前，先检查数据的外部选项处理方式；
外部选项的定义（包含什么、频率多少）本身是识别假设。

\subsection{序数结果上的社会偏好：Ambuehl--Bernheim（AER 2026）}

\notehead{模型}
替他人决策时，决策者知道他人对方案的\textbf{排序}但不知强度。
实验发现：多数人对“避免他人的最差结果”赋的权重，
高于“促成他人的最好结果”——排序信息被\textbf{位置性}使用
（类似 §4.6.2 最大最小倾向，但建立在序数信息上），
且该原则跨国稳定。

\notehead{可借鉴}
代议决策（家长为子女、受托人为受益人、政府为公民）
的福利函数不该是基数值的简单加总；
“防止最差”权重项（如 $\min$ 型或位置加权）更贴近真实社会偏好。

\subsection{模糊偏好的更新：Sadowski--Sarver（AER: Insights 2026）}

\notehead{模型与新质}
多先验/非期望效用（§\ref{subsec:meu}）在动态环境下有内在张力：
更新规则不能同时满足\textbf{动态一致}与\textbf{结果主义}
（只按事后偏好行事）。文章用演化方法论证：
动态一致是演化最优的硬要求，而结果主义的系统性违反反而最优——
为“事先承诺 vs 事后反悔”提供了演化基础。

\notehead{可借鉴}
模糊模型用于政策评估时，必须显式声明更新规则；
不同的更新规则下，同一政策的福利排序可以不同。

\subsection{异质个体的福利评估与有限注意：JPE 2025 两例}

\begin{itemize}
  \item \textbf{D\'avila \& Schaab}（JPE 2025）：
        个体在偏好与约束上多维异质时，福利变化分解为
        个体层面边际福利权重的加总——
        纠正“代表性消费者”对再分配效应的系统性遗漏；
  \item \textbf{Angeletos \& Sastry}（JPE 2025）：
        需求函数的斜率依赖\textbf{注意分配}：
        同一效用函数下，价格是否被注意决定需求弹性，
        名义刚性获得偏好以外的微观基础。
\end{itemize}

\subsection{中国问题中的效用函数：延迟退休与生育（详见 \S 4.6）}

\begin{center}
\footnotesize
\setlength{\tabcolsep}{4pt}
\begin{tabular}{p{1.9cm}p{3.0cm}p{4.6cm}p{4.6cm}}
\toprule
文献 & 模型 & 效用形式 & 关键结论 \\
\midrule
MDPI（2024） & 两期 OLG + 内生生育 & $\ln c_1+\beta\ln c_2+\phi\ln n+\eta\ln d$ & 延迟退休、缴费率与家庭赡养转移负相关 \\
SSRN（2024） & 多期 OLG + 祖辈照护 & 内生生育 + 隔代照料时间 & 延迟退休 5 年或使生育率降逾 25\% \\
IMF（2026） & 80 代 OLG、城乡分割 & 生命周期 CRRA 型 & 延迟退休 3 年：GDP +5.6\%（2050），养老金支出占比 15.3$\to$11.9\% \\
系统工程（2024） & OLG 内生生育与退休 & 生育补贴入预算约束 & 补贴促生育且促延迟退休；养老金呈倒 U \\
\bottomrule
\end{tabular}
\end{center}
共同特征：全部使用对数/CRRA 可加效用以换取闭式解与可校准性；
差异在\textbf{效用的自变量}——是否包含子女数量 $n$、父母消费 $d$、
祖辈照料时间。§4.9 完整复现此类模型的推导与计算。

\subsection{前沿速查表（2021--2026）}

\begin{center}
\small
\begin{tabular}{lll}
\toprule
方向 & 传统写法 & 前沿改写 \\
\midrule
福利稳健性 & 指定形式 $\to$ 算福利 & 形式无关的稳健半径（AER 2025） \\
偏好参数 & 结构常数 & 时变偏好面板（JPE 2025） \\
收入效应 & 位似（份额恒定） & 非位似、恩格尔斜率恒定（ECMA 2021） \\
选择架构 & 全合并、全理性 & 括定可检验（AER 2024） \\
数据基础 & 外部选项随手聚合 & 可观测性边界（AER 2026） \\
社会偏好 & 基数加总 & 序数位置加权（AER 2026） \\
模糊与动态 & 单一先验 & 多先验 + 更新规则声明（AERI 2026） \\
个体性 & 代表性消费者 & 多维异质 + 注意内生（JPE 2025） \\
\bottomrule
\end{tabular}
\end{center}

% ============================================================
\section{做研究时怎样选效用函数：十步流程}
% ============================================================

选错效用函数的后果有两类：\textbf{解不出}（形式太复杂）与
\textbf{解出来但结论是假的}（形式排除了真机制）。
下面的流程把“选函数”变成一组可执行的决定。每一步都给出
\textbf{操作}与\textbf{检查点}。

\subsection{第 1 步：界定问题的决策结构}

\notehead{操作}
回答三个是非题：(a) 选择是\textbf{一次性的还是序列的}（静态/动态）？
(b) 结果\textbf{确定还是不确定}（风险/不确定性）？
(c) 效用是否依赖\textbf{他人的选择或分配}（社会偏好/攀比）？
三个答案把候选集切到 §4.1--§4.5 的某一格。

\notehead{检查点}
延迟退休研究：序列（生命周期）+ 部分不确定（寿命、收入冲击）+ 含代际分配
（子女、父母）——落在“动态 + 风险 + 家庭内社会偏好”格，
直接排除静态单一商品形式。

\subsection{第 2 步：写清选择变量与状态变量}

\notehead{操作}
列两栏：个体\textbf{选择什么}（消费、储蓄、劳动、生育、退休时点），
\textbf{什么被外生给定或随状态演化}（工资、年龄、健康、政策参数、习惯存量）。
效用函数的自变量 $=$ 所有进入主观评价的变量，不能漏，也不必多。

\notehead{检查点}
漏掉关键自变量是高频错误：研究生育却不把子女数量/质量写入效用，
等于假设生育决策与偏好无关；研究退休却不写入闲暇，等于假设退休无收益。

\subsection{第 3 步：列出必须具备的偏好性质清单}

\notehead{操作}
对照第三章 §4.5 的翻译词典，写下本研究\textbf{不能排除}的性质：
\begin{center}
\small
\begin{tabular}{ll}
\toprule
研究情境 & 必须具备的性质（对应函数特征） \\
\midrule
需求份额随收入变化（必需品/奢侈品） & 非位似（基本需求型、非位似 CES） \\
单市场局部分析、要精确福利度量 & 无财富效应（拟线性） \\
宏观加总、要平衡增长路径 & 位似 + CRRA 跨期 \\
跨期与风险要分离 & 递归形式（EIS 与 RRA 解绑） \\
政策实验含“拖延/自控” & 动态不一致（$\beta$-$\delta$） \\
分配议题 & 差异厌恶项或攀比项 \\
\bottomrule
\end{tabular}
\end{center}

\notehead{检查点}
每条性质都应能回答：“如果形式不满足这条，我的核心机制会不会消失？”
会消失才写入清单——清单越短，形式越简单。

\subsection{第 4 步：可处理性分级}

\notehead{操作}
按计算资源选精度：
\begin{center}
\small
\setlength{\tabcolsep}{4pt}
\begin{tabular}{lll}
\toprule
级别 & 形式 & 代价 \\
\midrule
闭式解 & 对数、乘积型、CES、CARA-正态 & 弹性结构僵化 \\
解析条件 + 数值解 & CRRA、递归型、习惯形成 & 需校准与算法 \\
完全数值 & 非位似多部门、多维异质 & 检验成本高 \\
\bottomrule
\end{tabular}
\end{center}
经验法则：\textbf{理论命题用闭式解层级，政策定量用数值层级}，
两者用同一偏好族衔接（先对数后 CRRA）。

\subsection{第 5 步：参数从哪来——校准与识别}

\notehead{操作}
给每个参数指定来源：(i) 借用文献估计值（说明出处与适用性）；
(ii) 用矩条件校准（储蓄率、劳动份额、生育率等稳态矩）；
(iii) 结构估计（需求系统、选择数据）。
写下每个参数的\textbf{识别方程}：哪条数据矩决定哪个参数。

\notehead{检查点}
参数数 $>$ 独立矩数 $=$ 识别不足，必须砍参数（固定一部分）或换更简单形式。

\subsection{第 6 步：写出基准模型并求显式解}

\notehead{操作}
用最简可行形式完成全流程：一阶条件 $\to$ 闭式解或数值解
$\to$ 比较静态（每个政策参数对各选择变量的偏导，带符号判断）。
推导中每一步保留经济解释（如“$\partial n^*/\partial x$ = 收入效应 $-$ 照料成本效应”）。

\notehead{检查点}
若比较静态的符号无法从参数符号判断（含混），
说明机制的相对大小需要校准——这正是数值化的入口。

\subsection{第 7 步：对照数据做外部验证}

\notehead{操作}
用\textbf{未用于校准}的矩检验模型：预测的弹性、份额、政策响应
是否落在独立估计的区间内。

\subsection{第 8 步：逐步复杂化}

\notehead{操作}
一次只加一个成分（习惯 $\to$ 递归 $\to$ 异质性……），
每加一步重跑第 6 步，记录\textbf{结论的增量}。
增量为零的成分删掉：模型复杂度必须由结论需求证明。

\subsection{第 9 步：函数形式稳健性（顶刊标准）}

\notehead{操作}
按 Kang--Vasserman（§4.7.2）的精神执行三层稳健性：
\begin{enumerate}
  \item \textbf{参数层}：每个参数在置信区间内扰动，重算核心结论；
  \item \textbf{函数层}：换偏好族重做（对数 $\leftrightarrow$ CRRA；
        CES $\leftrightarrow$ 非位似 CES；指数贴现 $\leftrightarrow$ $\beta$-$\delta$）；
  \item \textbf{形式半径}：报告“偏离基准形式多远结论反转”的临界值。
\end{enumerate}

\subsection{第 10 步：与文献的对话定位}

\notehead{操作}
明确你的形式相对文献基准的\textbf{偏离点}：
改了哪个自变量、解绑了哪对参数、放松了哪条性质，
以及这个偏离为何对问题是必要的（不是为复杂而复杂）。

\subsection{决策树}

\begin{center}
\begin{tikzpicture}[
  font=\footnotesize,
  q/.style={draw=AxisBlue, rounded corners=2pt, align=center, inner sep=3pt, fill=MidBlue!10},
  a/.style={align=center, inner sep=2pt},
  arr/.style={-{Stealth[length=2.2mm]}, thick, color=AxisBlue},
  node distance=0.62cm and 0.8cm
]
\node[q] (q1) {动态？};
\node[q, below left=of q1, xshift=-0.9cm] (s1) {静态：\\ §4.1 谱系};
\node[q, below right=of q1, xshift=0.9cm] (q2) {有风险？};
\node[q, below left=of q2, xshift=-0.5cm] (d1) {纯动态：\\ 指数贴现起步};
\node[q, below right=of q2, xshift=0.5cm] (q3) {EIS 与 RRA\\ 需解绑？};
\node[a, below left=of q3, xshift=-0.4cm] (r1) {否：CRRA\\ 期望效用};
\node[a, below right=of q3, xshift=0.4cm] (r2) {是：\\ 递归形式};
\node[q, below=3.0cm of q1] (q4) {收入/财富效应\\ 是一阶的？};
\node[a, below left=of q4, xshift=-1.2cm] (f1) {否：拟线性\\ （部分均衡）};
\node[a, below right=of q4, xshift=1.2cm] (q5) {是：位似够吗？\\ 够：CES 族\\ 不够：非位似};
\draw[arr] (q1) -- node[above left] {否} (s1);
\draw[arr] (q1) -- node[above right] {是} (q2);
\draw[arr] (q2) -- node[above left] {否} (d1);
\draw[arr] (q2) -- node[above right] {是} (q3);
\draw[arr] (q3) -- (r1);
\draw[arr] (q3) -- (r2);
\draw[arr] (q4) -- (f1);
\draw[arr] (q4) -- (q5);
\draw[dashed, MidBlue] (s1) -- (q4);
\draw[dashed, MidBlue] (d1) -- (q4);
\end{tikzpicture}
\end{center}
\figexplan{决策树自“是否动态”起步：否则进入静态谱系；若是则再问有无风险——无风险取纯动态（指数贴现起步），有风险则按 EIS 与 RRA 是否需解绑选择 CRRA 期望效用或递归形式。两支最终由虚线汇入“收入/财富效应是否为一阶”的判断：否则用拟线性做部分均衡，是则按是否够位似选择 CES 族或非位似形式。}

\noindent 树只解决“从哪出发”；
第 6--9 步的\textbf{求解—验证—复杂化—稳健性}循环决定“走多远”。

% ============================================================
\section{政策应用：延迟退休、生育选择与效用函数的选取}
% ============================================================

本节以“中国延迟退休 + 生育政策”为完整案例，
执行 §4.8 十步流程的全部环节：选函数（论证）→ 建模型 → 逐步推导
→ 数值计算 → 政策对照 → 敏感性。所有算例数字精确到四位小数，可手工复算。

\subsection{第 1--3 步：问题界定与函数选择的论证}

\notehead{问题结构}
个体生命周期两期：成年工作期（选择消费 $c_1$、储蓄 $s$、生育数 $n$）、
老年期（消费 $c_2$；政策允许老年期的比例 $x$ 继续工作——$x$ 即\textbf{延迟退休强度}，
$x=0.2$ 对应 60 岁退休、$x=0.5$ 对应 65 岁退休，老年期长度标准化为 1 中的份额）。
政策参数：养老金缴费率 $\theta$、老年工作工资 $w_2$、养老金 $P$。

\notehead{效用函数的选择（对照 §4.8 第 3 步清单）}
\[
U=\ln c_1+\beta\ln c_2+\phi\ln n,
\]
逐项论证：
\begin{enumerate}
  \item \textbf{为什么跨期可加}：基准模型不需要习惯/递归结构；
        生命周期两期的储蓄—生育联合决策在可加形式下有闭式解（第 4 步）。
  \item \textbf{为什么 $\ln$ 而非一般 CRRA}：$\mathrm{EIS}=1$ 是中性基准；
        且 $\ln$ 使生育决策有\textbf{常数支出份额}结构，
        $n^*$ 闭式可解、比较静态符号清晰。$\gamma\ne1$ 的稳健性在 §\ref{subsec:crra-check} 检验。
  \item \textbf{为什么 $\phi\ln n$}：生育要成为\textbf{选择}，子女必须进入效用
        （第 2 步：不漏自变量）；$\ln n$ 意味着子女“越多越好但边际收益递减”，
        且养育支出占资源比例为常数 $\phi/(1+\phi+\beta)$——校准方便。
  \item \textbf{为什么暂不写父母消费/祖辈照料}：基准保持最小；
        照料渠道作为\textbf{机制扩展}在 §\ref{subsec:care} 引入（第 8 步：一次加一个成分）。
\end{enumerate}

\subsection{基准模型与完整推导}\label{subsec:baseline}

\notehead{预算约束}
成年期：工资 $w$（标准化为 $1$），缴费率 $\theta$，单位养育成本占工资 $v$：
\[
c_1=(1-\theta)w-s-vwn.
\]
老年期：比例 $x$ 工作（税后工资 $(1-\theta)w_2$），比例 $1-x$ 领取养老金 $P$，
储蓄按总回报 $R$ 积累：
\[
c_2=Rs+\underbrace{x(1-\theta)w_2+(1-x)P}_{\displaystyle\equiv\ \Omega(x)},
\]
其中 $\Omega(x)$ 是\textbf{老年期非资本收入}，
$\Omega'(x)=(1-\theta)w_2-P$：只要老年工作税后收入高于养老金，延迟退休就增加终身资源。

\notehead{一阶条件}
对 $s$：$-\dfrac{1}{c_1}+\beta\dfrac{R}{c_2}=0\ \Longrightarrow\ c_2=\beta Rc_1$。

对 $n$：$-\dfrac{vw}{c_1}+\dfrac{\phi}{n}=0\ \Longrightarrow\ n^*=\dfrac{\phi c_1}{vw}$。

\notehead{闭式解（逐步）}
\textbf{第一步}，把生育 FOC 代入成年预算消去 $n$：
$vwn=\phi c_1$，故
\[
c_1=(1-\theta)w-s-\phi c_1
\ \Longrightarrow\
c_1=\frac{(1-\theta)w-s}{1+\phi}.
\]
\textbf{第二步}，代入储蓄 FOC $c_2=\beta Rc_1$ 与老年预算 $c_2=Rs+\Omega$：
\[
\beta R\,\frac{(1-\theta)w-s}{1+\phi}=Rs+\Omega
\ \Longrightarrow\
\beta\bigl[(1-\theta)w-s\bigr]=(1+\phi)\Bigl(s+\frac{\Omega}{R}\Bigr),
\]
\[
\boxed{\ s^*=\frac{\beta(1-\theta)w-(1+\phi)\,\Omega/R}{1+\phi+\beta}\ }
\]
\textbf{第三步}，回代求 $c_1^*$：
\[
(1-\theta)w-s^*
=\frac{(1+\phi)\bigl[(1-\theta)w+\Omega/R\bigr]}{1+\phi+\beta}
\ \Longrightarrow\
\boxed{\ c_1^*=\frac{(1-\theta)w+\Omega/R}{1+\phi+\beta}\ }
\]
即\textbf{终身资源的现值}（劳动收入 $+$ 老年非资本收入的贴现）除以总权重 $1+\phi+\beta$。
\textbf{第四步}，其余变量：
\[
\boxed{\ c_2^*=\beta R\,\frac{(1-\theta)w+\Omega/R}{1+\phi+\beta},
\qquad
n^*=\frac{\phi}{vw}\cdot\frac{(1-\theta)w+\Omega/R}{1+\phi+\beta}\ }
\]

\notehead{比较静态（符号分析）}
\[
\frac{\partial n^*}{\partial x}
=\frac{\phi}{vw(1+\phi+\beta)}\cdot\frac{\Omega'(x)}{R}
=\frac{\phi\bigl[(1-\theta)w_2-P\bigr]}{vwR(1+\phi+\beta)}.
\]
基准模型中 $\Omega'(x)>0$（工作收入高于养老金），
故 $\partial n^*/\partial x>0$：\textbf{延迟退休提高终身收入，收入效应推升生育}；
同时 $\partial s^*/\partial x=-\dfrac{1+\phi}{1+\phi+\beta}\cdot\dfrac{\Omega'(x)}{R}<0$：
老年收入保障越好，储蓄动机越弱。
这是\textbf{只有收入效应}时的结论——机制不完整，见 §\ref{subsec:care}。

\subsection{数值算例（基准）}

取参数（25 年一期）：$w=1$，$w_2=1.2$，$\theta=0.16$，$P=0.4$，
$R=1.8$，$\beta=0.62$（$\approx0.98^{25}$），$\phi=0.35$，$v=0.12$。
则 $1+\phi+\beta=1.97$，$(1-\theta)w=0.84$，
$\Omega(x)=0.4+0.608x$。

\begin{center}
\small
\begin{tabular}{lccc}
\toprule
 & $x=0.2$（60 岁退） & $x=0.5$（65 岁退） & 变化 \\
\midrule
$\Omega(x)$ & $0.5216$ & $0.7040$ & $+35.0\%$ \\
$\Omega/R$ & $0.2898$ & $0.3911$ & \\
$c_1^*=(0.84+\Omega/R)/1.97$ & $0.5735$ & $0.6249$ & $+9.0\%$ \\
$c_2^*=\beta Rc_1^*$ & $0.6400$ & $0.6974$ & $+9.0\%$ \\
$s^*$ & $0.0658$ & $-0.0037$ & 储蓄动机消失 \\
$n^*=\phi c_1^*/(vw)$ & $1.6727$ & $1.8227$ & $+9.0\%$ \\
\bottomrule
\end{tabular}
\end{center}
（复算示例：$c_1^*(0.2)=(0.84+0.2898)/1.97=1.1298/1.97=0.5735$；
$n^*(0.2)=0.35\times0.5735/0.12=0.2007/0.12=1.6727$。）

\notehead{解读}
在只有收入效应的世界里，延迟退休 5 年使终身资源增加、
生育率从 $1.67$ 升至 $1.82$——与直觉和中国情境证据\textbf{相反}。
模型缺了关键一环：\textbf{祖辈照料}。这是函数形式（自变量集合）层面的问题，
不是参数问题。

\subsection{机制扩展：祖辈照料渠道}\label{subsec:care}

\notehead{修改}
中国家庭中祖辈是托育的主要供给者；延迟退休把祖父母的时间从照料中抽走。
把这一机制压缩为\textbf{有效养育成本随 $x$ 上升}：
\[
v(x)=v_0\bigl[1+\mu(x-0.2)\bigr],\qquad \mu>0,
\]
$x=0.2$ 为政策前基准，$\mu$ 衡量照料短缺的成本弹性。
解的形式不变（§\ref{subsec:baseline} 第四步），只是 $v$ 换成 $v(x)$：
\[
n^*(x)=\frac{\phi\bigl[(1-\theta)w+\Omega(x)/R\bigr]}{v(x)\,w\,(1+\phi+\beta)}
=\frac{\text{收入效应（分子，随 }x\text{ 升）}}{\text{照料成本效应（分母，随 }x\text{ 升）}}.
\]

\notehead{数值算例（含照料渠道）}
取 $\mu=1.5$：$v(0.5)=0.12\times(1+1.5\times0.3)=0.12\times1.45=0.174$。
\[
n^*(0.5)=\frac{0.35\times0.6249}{0.174}=\frac{0.2187}{0.174}=1.2571.
\]
\textbf{对照}：$n^*$ 从 $1.6727$ 降至 $1.2571$，降幅 $24.9\%$——
与近年多期 OLG 定量研究“延迟退休 5 年或使生育率降逾 25\%”的量级一致
（§4.7.9 表中 SSRN 2024）。

\begin{center}
\small
\begin{tabular}{lccc}
\toprule
延迟退休的生育效应 & 收入效应 & 照料成本效应 & 净效应 \\
\midrule
基准模型（无照料渠道） & $+9.0\%$ & --- & $+9.0\%$ \\
扩展模型（$\mu=1.5$） & $+9.0\%$ & $-31.0\%$ & $\mathbf{-24.9\%}$ \\
\bottomrule
\end{tabular}
\end{center}
（复算：$1.2571/1.6727-1=-0.2485$；照料渠道单独作用：
$1.8227\times0.12/0.174=1.2571$，即 $1/1.45-1=-31.0\%$。）

\notehead{方法论要点}
同一个效用函数 $U=\ln c_1+\beta\ln c_2+\phi\ln n$，
只是让\textbf{成本参数内生地依赖政策}，结论就翻转。
这说明：政策模型中“效用函数的选择”不仅是选曲率，
更是选\textbf{哪些价格/成本被政策移动}。

\subsection{生育数量约束：计划生育的效用分析}\label{subsec:quota}

\notehead{建模}
加入数量约束 $n\le\bar n$（$\bar n=1$ 对应独生子女政策、
$\bar n=2$ 对应两孩政策）。库恩—塔克条件：
\[
\frac{\phi}{n}-\frac{vw}{c_1}+\lambda=0,\qquad
\lambda\ge0,\quad \lambda(\bar n-n)=0.
\]
基准参数下自由解 $n^*=1.6727>1=\bar n$：约束\textbf{紧}，$n=\bar n=1$，
储蓄问题变为 $\max\ \ln c_1+\beta\ln c_2$（$n$ 固定）：
\[
s^{**}=\frac{\beta\bigl[(1-\theta)w-vw\bigr]-\Omega/R}{1+\beta}
=\frac{0.62\times0.72-0.2898}{1.62}=\frac{0.1566}{1.62}=0.0967,
\]
$c_1^{**}=0.72-0.0967=0.6233$，$c_2^{**}=\beta Rc_1^{**}=0.6956$。
\textbf{影子价格}（放松一个生育配额的边际效用价值）：
\[
\lambda=\frac{\phi}{\bar n}-\frac{vw}{c_1^{**}}
=0.35-\frac{0.12}{0.6233}=0.35-0.1925=0.1575>0\quad(\text{约束确紧}).
\]

\notehead{福利损失的计算}
\[
\begin{aligned}
U_{\text{自由}}&=\ln 0.5735+0.62\ln 0.6400+0.35\ln 1.6727\\
&=-0.5561+0.62\times(-0.4463)+0.35\times0.5141
=-0.5561-0.2767+0.1799=-0.6529,\\
U_{\text{配额}}&=\ln 0.6233+0.62\ln 0.6956+0.35\ln 1
=-0.4727+0.62\times(-0.3631)+0=-0.6979.
\end{aligned}
\]
$\Delta U=-0.6979+0.6529=-0.0450$。
换算为消费等价变化：$e^{\Delta U}-1=e^{-0.0450}-1=-4.4\%$——
独生子女配额等价于\textbf{终身消费下降约 4.4\%} 的福利损失（基准参数下）。

\notehead{政策放开的预测}
$\bar n$ 提高到 $2$：自由解 $1.6727<2$，约束不再紧，
生育率升至 $1.67$ 而非 $2$。
\textbf{这正解释了中国全面两孩后的现实}：配额放开后生育率远低于 $2$——
约束早已从“数量约束”变为“经济约束”（$v$ 高、$w$ 增长放缓、$\Omega$ 预期）。
效用框架给出可检验的判别式：约束紧不紧，看影子价格
$\lambda=\phi/\bar n-vw/c_1$ 的符号。

\subsection{扩展：数量—质量权衡与内生退休}

\notehead{数量—质量}
$U=\ln c_1+\beta\ln c_2+\phi\ln n+\psi\ln q$，
养育成本 $n(vw+q)$（$q$ 为每个子女的人力资本投入）。
一阶条件给出
\[
\frac{\phi/n}{vw+q}=\frac{\psi/q}{n}
\ \Longrightarrow\
\phi q=\psi n(vw+q)
\ \Longrightarrow\
q^*=\frac{\psi nvw}{\phi-\psi n}\quad(\text{要求 }\phi>\psi n),
\]
$q^*$ 是 $n$ 的增函数方向反过来成立：外生压低的 $n$（配额）
\textbf{推高}每孩投入 $q$——独生子女政策与“精细化育儿/教育军备竞赛”
的关联在此结构中是一个内生结果。

\notehead{内生退休}
让 $x$ 成为选择变量，老年期效用加入闲暇：
$\ln c_2+\chi\ln(1-x)$。一阶条件
\[
\frac{\chi}{1-x}=\frac{\Omega'(x)}{c_2}
\ \Longrightarrow\
1-x^*=\frac{\chi c_2}{(1-\theta)w_2-P}.
\]
取 $\chi=0.5$：$1-x^*=0.5\times0.64/0.608=0.526$，$x^*\approx0.47$——
个体自愿的退休强度与政策 $x=0.5$ 接近，说明当前政策接近
\textbf{紧约束边界}：小幅提高养老金 $P$ 就会使自愿退休早于法定，
政策实效将依赖劳动市场（老年就业可得性）而非文件规定。

\subsection{敏感性：把对数换成 CRRA（曲率取 2）}\label{subsec:crra-check}

效用 $\dfrac{c_1^{1-\gamma}}{1-\gamma}+\beta\dfrac{c_2^{1-\gamma}}{1-\gamma}+\phi\ln n$，
$\gamma=2$。一阶条件变为
\[
c_2=(\beta R)^{1/\gamma}c_1=1.0564\,c_1,\qquad
n^*=\frac{\phi\,c_1^{\gamma}}{vw}.
\]
代入预算（$x=0.2$，同参数），得方程
$0.35c_1^2+1.5869c_1-1.1298=0$，正根 $c_1^*=0.6256$；
\[
n^*=\frac{0.35\times0.6256^2}{0.12}=1.1416,\qquad
s^*=0.0774,\qquad c_2^*=0.6609.
\]
\textbf{对照}：$\gamma=1$（对数）时 $n^*=1.6727$；$\gamma=2$ 时 $n^*=1.1416$（$-32\%$）。
曲率越大，消费平滑动机越强，以消费计价的养育成本显得越贵，
生育率越低。\textbf{政策结论的方向不变（配额紧、照料渠道主导），
但定量幅度对折现曲率高度敏感}——这正是 §4.8 第 9 步要求
报告“函数层稳健性”的原因。

\subsection{从算例到论文：研究路线}

\begin{enumerate}
  \item \textbf{基准}：本节两期模型（闭式解，机制透明）；
  \item \textbf{定量化}：扩展为多期 OLG（80 代），效用仍取对数可加族，
        校准目标：储蓄率、劳动参与、总和生育率、养老金替代率
        （参照 §4.7.9 表中 IMF 2026 的城乡分割结构）；
  \item \textbf{数据}：微观层面用 CHARLS/CFPS 估计 $\phi$、$\mu$（祖辈照料的成本弹性）、
        $\beta$；宏观层面校准 $\theta$、$P$、$R$；
  \item \textbf{反事实}：延迟退休 $\times$ 生育补贴的组合政策；
        生育补贴在模型中即降低 $v$——与延迟退休提高 $v(x)$ 对冲，
        存在使 $n^*$ 不变的\textbf{政策组合线}；
  \item \textbf{稳健性}：CRRA 曲率、习惯形成（老年消费的习惯存量）、
        递归偏好（寿命不确定性下 EIS 与 RRA 解绑）三个方向各做一遍。
\end{enumerate}

\notehead{本章收束}
效用函数不是“越真越好”，而是\textbf{“恰好包含所研究机制的最简形式”最好}。
本节的演示：同一个对数可加效用，(i) 在纯收入效应下给出错误符号的预测，
(ii) 加入照料成本内生化后给出与定量文献一致的幅度，
(iii) 加入数量约束后解释政策放开后的低生育，
(iv) 换曲率后给出稳健性区间——
选函数的每一步都是\textbf{可被数据检验、被文献对话}的决定。

\clearpage
\section{拓展：效用函数选择的识别、估计与稳健性}

\subsection{从理论形式到可估计对象}
写下效用函数只是研究设计的起点。一个可估计模型还必须明确参数、观测变量、误差结构与识别来源。以离散选择为例，随机效用写为
\[
U_{ij}=V_{ij}(\theta)+\varepsilon_{ij},
\]
其中 $V$ 是可解释部分，$\varepsilon$ 是未观测冲击。若误差服从独立同分布极值分布，则选择概率为 Logit 形式；若误差相关，必须改用 Nested Logit 或混合 Logit。

\begin{figure}[H]\centering
\begin{tikzpicture}[font=\scriptsize,>=stealth,scale=0.92,transform shape]
\node[draw,rounded corners,fill=gray!10,minimum width=2.3cm,minimum height=.8cm] (t) at (0,0) {理论偏好};
\node[draw,rounded corners,fill=gray!18,minimum width=2.3cm,minimum height=.8cm] (f) at (3.5,0) {函数形式};
\node[draw,rounded corners,fill=gray!26,minimum width=2.3cm,minimum height=.8cm] (d) at (7,0) {数据与变异};
\node[draw,rounded corners,fill=gray!34,minimum width=2.3cm,minimum height=.8cm] (e) at (10.5,0) {估计与检验};
\draw[->,thick] (t)--(f);\draw[->,thick] (f)--(d);\draw[->,thick] (d)--(e);
\node[align=center] at (5.25,-1.25) {每一步都可能引入不可识别参数\\不能用更复杂的函数弥补缺少的变异};
\end{tikzpicture}
\caption{效用函数研究的识别链条}
\end{figure}
\figexplan{四个方框依次表示理论、函数、数据和估计。箭头不是单纯的计算流程，而是识别约束逐步收紧的过程；若数据没有足够变异，形式越复杂并不意味着结论越可信。}

\subsection{尺度归一化与参数识别}
在随机效用模型中，若同时把系统效用与误差项乘以正数，选择概率不变。因此必须固定误差尺度，例如令误差方差为 $\pi^2/6$，或固定某一个价格系数。若不归一化，参数向量只在比例类意义下识别。

\begin{proposition}[Logit 概率]
若 $\varepsilon_{ij}$ 独立服从标准极值分布，则
\[
P_{ij}=\frac{\exp(V_{ij})}{\sum_k\exp(V_{ik})}.
\]
对任意常数 $c$，$P_{ij}(V_i+c)=P_{ij}(V_i)$；这体现了效用水平零点不可识别。
\end{proposition}

\subsection{弹性、福利与反事实}
价格弹性必须说明保持什么不变。马歇尔需求的自身价格弹性为
$\varepsilon_{ii}=\partial\log x_i/\partial\log p_i$；补偿弹性则保持效用不变。福利分析中，消费者剩余近似适用于小变化，但当收入效应显著时，应使用等价变差或补偿变差：
\[
EV=e(p^0,u^1)-e(p^0,u^0),\qquad
CV=e(p^1,u^1)-e(p^1,u^0).
\]
这两个量分别把福利变化放在初始价格和最终价格下衡量。

\begin{figure}[H]\centering
\begin{tikzpicture}[x=0.9cm,y=0.85cm,font=\scriptsize,scale=0.92,transform shape]
\draw[->] (0,0)--(6,0) node[right] {$q$};\draw[->] (0,0)--(0,4.5) node[above] {$p$};
\draw[thick] (0.8,4)--(5.4,.8) node[right] {$D$};
\draw[dashed] (0,2.7)--(5.4,2.7) node[right] {$p^0$};
\draw[dotted] (0,1.8)--(5.4,1.8) node[right] {$p^1$};
\fill[gray!25] (0,1.8)--(0,2.7)--(2.6,2.7)--cycle;
\node[align=center,fill=white] at (1.15,2.3) {福利变化\\面积};
\end{tikzpicture}
\caption{价格变化与福利面积}
\end{figure}
\figexplan{横轴为数量、纵轴为价格。两条水平线表示政策前后的价格，需求曲线与两线之间的灰色区域是小幅变化下的消费者剩余变化；大幅变化应回到支出函数定义。}

\subsection{稳健性：函数形式不是结论本身}
至少应报告四类稳健性：替换函数形式、替换误差分布、改变样本窗口、改变关键变量定义。若核心结论只在一种形式下成立，应将其表述为模型依赖结果，而非普遍事实。

\begin{example}[三种函数形式的比较]
同一政策问题可分别使用 Cobb--Douglas、CES 与 Translog。Cobb--Douglas 施加单位替代弹性，CES 允许恒定但不必为一的替代弹性，Translog 则允许弹性随消费束变化。比较三者的拟合、外推和参数稳定性，才能判断简洁性是否以错误限制为代价。
\end{example}

\subsection{研究清单}
\begin{enumerate}
\item 明确效用是序数、基数还是随机效用指数；
\item 列出每个参数的经济含义和识别来源；
\item 检查角点、边界与不可微点；
\item 说明福利度量采用 EV、CV 还是消费者剩余近似；
\item 报告至少一种替代函数与一种替代估计策略。
\end{enumerate}

\subsection{社会福利函数的精细校对}

\begin{proposition}[CES 社会福利的两个极限]
令
\[
W_\varepsilon(x)=\left(\frac1N\sum_{i=1}^Nx_i^{1-\varepsilon}\right)^{1/(1-\varepsilon)}.
\]
当 $\varepsilon\to0$ 时，$W_\varepsilon$ 收敛于算术平均；当 $\varepsilon\to1$ 时，收敛于几何平均；当 $\varepsilon\to\infty$ 时，收敛于 $\min_i x_i$。
\end{proposition}
\begin{proof}
第一极限直接代入 $1-\varepsilon=1$。令 $r=1-\varepsilon$，则 $\ln W_r=\frac1r\ln(\frac1N\sum_i x_i^r)$。在 $r\to0$ 时，$x_i^r=e^{r\ln x_i}=1+r\ln x_i+o(r)$，故
$\ln W_r\to N^{-1}\sum_i\ln x_i$，得到几何平均。最后，幂均值的最大项控制 $r\to-\infty$ 的极限，取对数后最大项对应原变量中的最小值，故极限为最小值。
\end{proof}

\begin{example}[同一可能性前沿上的切点]
在图形中，可能性前沿必须保持不变，变化的只能是社会福利等值线。功利主义等值线斜率为 $-1$；最大最小准则的有效等值线是水平线或垂直线的分段；CES 等值线在两者之间平滑变化。若分别为三种准则画出不同的前沿，会把“技术约束变化”误读为“评价准则变化”。
\end{example}

\begin{remark}[图 4.3 的读图规则]
图 4.3 中三幅面板共享坐标轴、共享可能性前沿和共享效用单位。左图的切点由 $W=u_1+u_2$ 的等值线决定，中图由最差者水平线决定，右图由 CES 等值线决定。图内标签统一放在白底节点中，避免压住前沿；箭头只标示福利改善方向，不承担曲线定义。
\end{remark}

\subsection{福利变化的计算核对}
对收入向量 $(10,2)$，算术平均为 $6$，几何平均为 $\sqrt{20}$。当 $\varepsilon=2$ 时，CES 福利为
\[
\left[\frac12(10^{-1}+2^{-1})\right]^{-1}=\frac1{0.3}=3.3333.
\]
因此不平等指数 $A=1-W/6=0.4444$。这里的 $A$ 是相对于同一算术平均的归一化指标，不能与使用不同收入单位或不同基准均值的指数直接比较。


\end{document}
